\documentclass[11pt]{article}

\usepackage[a4paper,margin=1in]{geometry}
\usepackage{amsmath,amssymb,amsthm,mathtools}
\usepackage[hidelinks]{hyperref}
\hypersetup{
  pdfauthor={Bangyen Pham},
  pdftitle={Coefficient Order Statistics of Multiples of a Power Below the Root 2},
  pdfsubject={The rows of polynomial multiples of (x-r)^L for 1<r<2: where the prefix description holds and where it fails},
  pdfkeywords={polynomial multiples, coefficient order statistics, reduced height, linear programming duality, computer-assisted proof},
  pdflang={en-US}
}

\newtheorem{theorem}{Theorem}[section]
\newtheorem{lemma}[theorem]{Lemma}
\newtheorem{corollary}[theorem]{Corollary}
\newtheorem{proposition}[theorem]{Proposition}
\newtheorem*{mainthm}{Main Theorem}
\newcommand{\R}{\mathbb R}
\newcommand{\N}{\mathbb N}
\DeclareMathOperator{\sgn}{sgn}

\title{Coefficient Order Statistics of Multiples of a Power Below the Root $2$}
\author{\shortstack{Bangyen Pham\\
\small Independent Researcher\\
\small\href{mailto:bangyenp@gmail.com}{\texttt{bangyenp@gmail.com}}\\
\small\href{https://orcid.org/0009-0006-9928-2088}{ORCID: 0009-0006-9928-2088}}}
\date{September 2026}

\begin{document}
\maketitle

\begin{abstract}
For $r>1$ let $V_r(L,k)$ be the infimum, over monic real multiples $F$ of
$(x-r)^L$, of the $k$-th largest magnitude of a nonleading coefficient of
$F$.  At every root $r\ge2$ exempting the $k-1$ largest coefficients costs
exactly $k-1$ roots, and the worst exempted positions are the top ones:
$V_r(L,k)=\min_{i\le k}1/\nu_i(L-k+1)$, where $\nu_i(n)$ is a
prefix-constrained $\ell^1$ problem.  We show that for $1<r<2$ this prefix
identity holds on every diagonal $L-k=\mathrm{const}$ in all but finitely
many rows, and in every row when the diagonal's top row is at most one,
but fails in general: it fails in some row for every
$r\in[\frac{21}{20},\frac{19}{10}]$, and in the second row for every
$r\in[\frac{21}{20},\frac{139}{100}]$.  The tool is a finite reduction:
every row $k\ge2$ is the minimum of the preceding row on its diagonal and
finitely many terms, each attained by a polynomial with integer zeros whose
optimality, for rational $r$, is a finite exact check.
\end{abstract}

\noindent\textbf{Keywords.} Polynomial multiples; coefficient order
statistics; reduced height; linear programming duality; computer-assisted
proof.

\noindent\textbf{2020 Mathematics Subject Classification.} Primary 11C08; Secondary 26C05, 90C05.

\section{Introduction}\label{sec:intro}

For a nonzero real polynomial $F$ of degree $D$ let $b_k(F)$ be the
$k$-th entry of its nonleading coefficient magnitudes $|f_j|$, $j<D$, in
decreasing order with multiplicity.  Fix a real $r>1$ and put
\[
  V_r(L,k)=\inf\{b_k(F):\ F\text{ a monic real multiple of }(x-r)^L\},
  \qquad \beta_r(n)=V_r(n,1),
\]
the row $k$ of $(x-r)^L$ and the top row of $(x-r)^n$.  Exempting the
$k-1$ largest coefficients costs at least $k-1$ roots,
$V_r(L,k)\le\beta_r(L-k+1)$, and the companion
paper~\cite{coefficient-mass-rows} proves equality at every $r\ge2$
\cite[Theorem~5.3]{coefficient-mass-rows}.  This paper asks what happens
below the root $2$, where the last row
$V_r(L,L)=(r-1)^L<\beta_r(1)$ already shows that the rows are not all top
rows \cite[Corollary~2.4]{coefficient-mass-rows}.  The top row
$\beta_r(n)$ is the reduced height of $(x-r)^n$ in the sense of Dubickas
and Jankauskas~\cite{dj}, and the rows $k\ge2$ are the order-statistic
variant of~\cite{coefficient-mass}.

\paragraph{Setup.}
We use the notation of~\cite{coefficient-mass-rows}.  Put $a=1/(r-1)$ and,
for a finite $Z\subset\N_{>0}$, $r_Z(s)=\prod_{z\in Z}(1-s/z)$.  For real
polynomials $q$ put $\Phi_r(q)=\sum_{s\ge1}|q(s)|r^{-s}$ and
$\Phi_r(Z)=\Phi_r(r_Z)$, and for finite $S\subset\N_{>0}$ and $L>|S|$ let
$\mu_r(S,L)$ be the infimum of $\Phi_r(q)$ over real $q$ with
$\deg q\le L-1$, $q(0)=1$ and $q|_S=0$.  By linear-programming duality
\cite[Proposition~2.1]{coefficient-mass-rows}, $V_r(L,k)$ is the infimum
of $1/\mu_r(S,L)$ over exempted sets $|S|=k-1$ of backward distances, the
infimum is attained, and $\beta_r(n)=1/\mu_r(\varnothing,n)$.  The prefix
values are $\nu_i(n)=\mu_r([1,i-1],n+i-1)$, so $\nu_1(n)=1/\beta_r(n)$,
and with $n=L-k+1$ they bracket every row
\cite[Theorem~3.3]{coefficient-mass-rows}:
\[
  \frac1{\Omega_k(n)}\ \le\ V_r(L,k)\ \le\ \min_{1\le i\le k}\frac1{\nu_i(n)},
\]
where $\Omega_k(n)$ is the least value of $\sum_sW_k(s)r^{-s}|p(s)|$ over
real $p$ with $\deg p<n$ and $p(0)=1$, and $W_k(s)=\max_{i\le k}\binom{s-1}{i-1}$;
the loss $E_k(r)$ of that theorem comes from the first $2k-2$ positions.
For finite $Z$ and $c\ge1$, $\mathrm{ins}(Z,c)$ is $Z$ with $c$ inserted
and the elements from $c$ on moved up by one.
They also satisfy \cite[Section~5]{coefficient-mass-rows}
\begin{equation}\label{eq:prefixup}
  \nu_{i+1}(n)\le a\,\nu_i(n)\qquad(r>1,\ i,n\ge1).
\end{equation}
We say that the \emph{prefix identity} holds at $(r,L,k)$ if the upper
bound is attained:
\begin{equation}\label{eq:prefixid}
  V_r(L,k)=\min_{1\le i\le k}\frac1{\nu_i(n)}\qquad(n=L-k+1),
\end{equation}
so that the prefix $[1,i-1]$ is a worst exempted set for some $i\le k$.
At $r\ge2$ it holds everywhere, with the minimum at $i=1$.

\paragraph{Where the identity holds.}
For $1<r<2$ the identity holds for every $k$ whenever $\beta_r(n)\le1$
(Theorem~\ref{thm:onepolyrows}), and on every diagonal in all but finitely
many rows (Corollary~\ref{cor:longdiagfinite}): it holds whenever
$k\ge\max(K,\alpha(n-1))$, where $\alpha\ge1$ is any constant with
$\rho(\alpha)<1$ for an explicit function $\rho$ depending on $r$, and $K$
depends on $r$ and $\alpha$ (Theorem~\ref{thm:longdiag},
Proposition~\ref{prop:longdiagexplicit}).  In both regimes the minimum is
at $i=k$: the prefix $[1,k-1]$ is a worst exempted set.

\paragraph{Where it fails.}
The second row $V_r(n+1,2)$ is the minimum of the top row $\beta_r(n)$ and
finitely many exempted-position terms (Theorem~\ref{thm:secondrow}), and
every row $k\ge2$ is the minimum of $V_r(L-1,k-1)$ and finitely many
exempted-set terms (Theorem~\ref{thm:finiterows}).  Each term is attained
by a polynomial with integer zeros, and its optimality is a finite exact
check (Lemma~\ref{lem:vertexopt}).  With these reductions:
\begin{mainthm}
For every $r\in[\frac{21}{20},\frac{19}{10}]$ the prefix identity fails
in some row of some diagonal, and for every
$r\in[\frac{21}{20},\frac{139}{100}]$ it fails in the second row
(Theorems~\ref{thm:failinterval} and~\ref{thm:blockinterval}).
\end{mainthm}
Since the values $\mu_r(S,L)$ decrease with $r$
(Lemma~\ref{lem:rootmono}), a certificate at one root and upper bounds at
a smaller one give a failure at every root between them, and finitely many
certificates cover each interval.  Above $\frac{139}{100}$ the failures are
in higher rows: the term of the exempted block $[2,k]$ is at least a prefix
value of a longer prefix and a smaller degree, with equality when its
minimizer vanishes at $1$ (Lemma~\ref{lem:block}), and the certified rows
grow as $r\to2$, up to $k=384$.  Exact examples at single roots show that
the identity can fail in the second row even when $\nu_2(n)>\nu_1(n)$
(Proposition~\ref{prop:secondrowhyp}), and can fail in the third row of a
diagonal while holding in the second (Propositions~\ref{prop:secondrowexact}
and~\ref{prop:thirdrowexact}).

The proofs of these three propositions and of
Theorems~\ref{thm:failinterval} and~\ref{thm:blockinterval} are
computer-assisted: they list the certifying zero sets and reduce the
claims to finitely many comparisons of explicit rationals, which the
accompanying tests carry out in exact arithmetic.  Every other result is
proved by hand.

\paragraph{Open.}
We do not know whether the identity fails for every $1<r<2$, that is,
also below $\frac{21}{20}$ and above $\frac{19}{10}$, nor whether the
least failing row tends to infinity as $r\to2$.

\paragraph{Organization.}
Section~\ref{sec:onepoly} proves the identity when $\beta_r(n)\le1$ by
using one polynomial for all exempted sets, and Section~\ref{sec:longdiag}
proves it on long diagonals.  Section~\ref{sec:secondrow} reduces the
second row to finitely many terms and certifies its failures, and
Section~\ref{sec:finiterows} does the same for every row.

\paragraph{Conventions.}
All logarithms are natural, $\N_{>0}=\{1,2,\ldots\}$, and $[u,v]$ denotes
the integers in that range.  Results of the companion papers are cited by
their numbers there.  Table~\ref{tab:notation} lists the notation used in
more than one place.

\begin{table}[tp]
\centering
\small
\renewcommand{\arraystretch}{1.15}
\begin{tabular}{@{}l p{0.47\textwidth} l@{}}
\hline
Symbol & Meaning & Defined in\\
\hline
$V_r(L,k)$, $\beta_r(n)$ & row $k$ of $(x-r)^L$; top row $V_r(n,1)$
  & Section~\ref{sec:intro}\\
$a$, $r_Z(s)$ & $1/(r-1)$; $\prod_{z\in Z}(1-s/z)$ & Section~\ref{sec:intro}\\
$\Phi_r(q)$, $\mu_r(S,L)$ & $\sum_{s\ge1}|q(s)|r^{-s}$; its infimum over
  $\deg q<L$, $q(0)=1$, $q|_S=0$ & Section~\ref{sec:intro}\\
$\nu_i(n)$ & prefix value $\mu_r([1,i-1],n+i-1)$ & Section~\ref{sec:intro}\\
prefix identity & $V_r(L,k)=\min_{i\le k}1/\nu_i(n)$ & \eqref{eq:prefixid}\\
$\omega_i(s)$ & $\binom{s-1}{i-1}r^{-s}$ & Section~\ref{sec:onepoly}\\
$\mathcal Q_n$ & real $p$ with $\deg p\le n-1$ and $p(0)=1$
  & Section~\ref{sec:onepoly}\\
$A_i(p)$ & $\sum_s\omega_i(s)|p(s)|=\Phi_r(r_{[1,i-1]}p)$
  & Section~\ref{sec:onepoly}\\
$\Lambda_k(n)$ & $\inf_{p\in\mathcal Q_n}\max_{i\le k}A_i(p)$
  & Theorem~\ref{thm:onepoly}\\
$w_k$ & $\omega_k-\omega_{k-1}$ & Section~\ref{sec:longdiag}\\
$\tau_k(n)$ & $\inf_{p\in\mathcal Q_n}\sum_{s\ge2k-2}w_k(s)|p(s)|$
  & Section~\ref{sec:longdiag}\\
$N_k(r)$ & largest excess of an $\omega_i$, $i<k$, over $\omega_k$ on
  $[1,2k-3]$ & Section~\ref{sec:longdiag}\\
$\Xi(k)$, $\rho(\alpha)$ & explicit bounds for long diagonals
  & Proposition~\ref{prop:longdiagexplicit}\\
$\eta_\sigma(n)$ & $\mu_r(\{\sigma\},n+1)$, one exempted position
  & Section~\ref{sec:secondrow}\\
$\Psi_r(p)$, $R_p(\sigma)$ & $\sum_ss|p(s)|r^{-s}$;
  $\sum_{s>\sigma}(s-\sigma)|p(s)|r^{-s}$ & Section~\ref{sec:secondrow}\\
$e_y$, $G_y$ & the vertex test & Lemma~\ref{lem:vertexopt}\\
$p_*$, $\sigma_*$ & a minimizer for $\nu_1(n)$; its tail threshold
  & Theorem~\ref{thm:secondrow}\\
$Y_n$, $Y'_n$, $Z_S(n)$ & named zero sets of the exact certificates
  & Section~\ref{sec:secondrow}\\
$\phi_m(s,t)$, $\Delta_{q,m}(t)$ & weight and gap for $m$ far zeros
  & Section~\ref{sec:finiterows}\\
$t_m(q)$ & least integer $t\ge1$ with $\Delta_{q,m}(t)\le0$
  & Lemma~\ref{lem:farzeros}\\
$T(N)$, $\mathcal S$ & thresholds and the finite tree of exempted sets
  & Theorem~\ref{thm:finiterows}\\
\hline
\end{tabular}
\caption{Notation.}\label{tab:notation}
\end{table}

\section{One polynomial for all exempted sets}\label{sec:onepoly}

This section proves~\eqref{eq:prefixid} for $1<r<2$ whenever $\beta_r(n)\le1$,
for every $k$ (Theorem~\ref{thm:onepolyrows}).  There the minimum is at
$i=k$: the prefix $[1,k-1]$ is a worst exempted set.

It rests on two facts: one polynomial $p$ serves all exempted sets at
once, at the cost of the largest of its prefix sums
(Theorem~\ref{thm:onepoly}); and these prefix sums grow at least by the
factor $\nu_1(n)=1/\beta_r(n)$ from one index to the next
(Lemma~\ref{lem:prefixshift}).  For the first,
\cite[Lemma~4.1(b)]{coefficient-mass-rows} pushes every exempted position down into an
initial block while the fixed factor $p$ rides along.  Nothing is inserted,
and the only optimality used is that of $p$ for the last index.

Put $\omega_i(s)=\binom{s-1}{i-1}r^{-s}$ for $i\ge1$ (zero for $s<i$),
let $\mathcal Q_n$ be the set of real polynomials $p$ with
$\deg p\le n-1$ and $p(0)=1$, and put
\[
  A_i(p)=\sum_{s\ge1}\omega_i(s)\,|p(s)|=\Phi_r\bigl(r_{[1,i-1]}\,p\bigr)
  \qquad(i\ge1),
\]
for real polynomials $p$, so that $\nu_i(n)=\inf_{p\in\mathcal Q_n}A_i(p)$.
The second expression holds because $|r_{[1,i-1]}(s)|=\binom{s-1}{i-1}$
for $s\ge1$.  
By \cite[Lemma~4.4]{coefficient-mass-rows}, which rests on the convexity
of \cite[Lemma~4.1(b)]{coefficient-mass-rows}, every real polynomial $p$
and finite $S\subset\N_{>0}$ satisfy
\begin{equation}\label{eq:prefixpush}
  \Phi_r(r_S\,p)\le\max_{1\le i\le|S|+1}A_i(p).
\end{equation}

\begin{theorem}[One polynomial for all exempted sets]\label{thm:onepoly}
Let $r>1$, $1\le k\le L$, $n=L-k+1$ and $p\in\mathcal Q_n$.  Every monic
multiple $F$ of $(x-r)^L$ has
\[
  b_k(F)\ \ge\ \frac1{\max_{1\le i\le k}A_i(p)} .
\]
Hence, with $\Lambda_k(n)=\inf_{p\in\mathcal Q_n}\max_{1\le i\le k}A_i(p)$,
\[
  \frac1{\Omega_k(n)}\ \le\ \frac1{\Lambda_k(n)}\ \le\ V_r(L,k)\ \le\
  \min_{1\le i\le k}\frac1{\nu_i(n)} .
\]
\end{theorem}

\begin{proof}
Let $S$ be the exempted set of \cite[Proposition~2.1(a)]{coefficient-mass-rows}, so
$|S|=k-1$.  The polynomial $q=r_Sp$ has $\deg q\le k-1+n-1=L-1$, $q(0)=1$
and $q|_S=0$.  So $b_k(F)\,\Phi_r(q)\ge1$, and~\eqref{eq:prefixpush}
bounds $\Phi_r(q)$.  The upper bound is \cite[Theorem~3.3]{coefficient-mass-rows}.  For
the first inequality, $W_k(s)=\max_{1\le i\le k}\binom{s-1}{i-1}$ in the
notation of \cite[Theorem~3.3]{coefficient-mass-rows}, so
$\sum_sW_k(s)r^{-s}|p(s)|\ge\max_{i\le k}A_i(p)$ for every $p$.
\end{proof}

The theorem moves the maximum over $i$ out of the sum in
\cite[Theorem~3.3]{coefficient-mass-rows}, and this loses nothing when one $p$ is nearly
optimal for all prefix problems at once.  The prefix identity~\eqref{eq:prefixid} holds at $(r,L,k)$ whenever
$\Lambda_k(n)=\max_{i\le k}\nu_i(n)$.  For $n=1$ we have
$\mathcal Q_1=\{1\}$ and $A_i(1)=a^i$.  So the theorem gives
$V_r(L,L)\ge r-1$ for $r\ge2$ and $V_r(L,L)\ge(r-1)^L$ for $r\le2$, the
lower bounds of \cite[Corollary~2.4]{coefficient-mass-rows}; its proof is this argument,
through \cite[Lemma~4.4]{coefficient-mass-rows} with $p=1$.

\begin{lemma}[Shifting the prefix index]\label{lem:prefixshift}
Let $r>1$, $n\ge1$, $i,m\ge1$, and let $p$ be real with $\deg p\le n-1$.
Then
\[
  A_{i+m}(p)\ \ge\ \nu_m(n)\,A_i(p).
\]
Consequently $\nu_{i+m}(n)\ge\nu_i(n)\,\nu_m(n)$, and
$\nu_i(n)\ge\nu_1(n)^i=\beta_r(n)^{-i}$.
\end{lemma}

\begin{proof}
Counting the $(i+m-1)$-element subsets of $[1,s-1]$ by their $i$-th
smallest element $u$ gives
$\sum_{u=1}^{s-1}\binom{u-1}{i-1}\binom{s-u-1}{m-1}=\binom{s-1}{i+m-1}$, that is,
$\omega_{i+m}(s)=\sum_{u+t=s,\ u,t\ge1}\omega_i(u)\,\omega_m(t)$.  All terms
being nonnegative,
\[
  A_{i+m}(p)=\sum_{u\ge1}\omega_i(u)\sum_{t\ge1}\omega_m(t)\,|p(u+t)| .
\]
If $p(u)\ne0$, the polynomial $t\mapsto p(u+t)/p(u)$ lies in $\mathcal Q_n$,
so the inner sum is at least $\nu_m(n)|p(u)|$; if $p(u)=0$ this is trivial.
This proves the inequality.  Taking the infimum over $p\in\mathcal Q_n$
gives $\nu_{i+m}(n)\ge\nu_m(n)\nu_i(n)$, and $\nu_1(n)=1/\beta_r(n)$.
\end{proof}

Together with~\eqref{eq:prefixup},
$\nu_1(n)\,\nu_i(n)\le\nu_{i+1}(n)\le a\,\nu_i(n)$, and both inequalities are
equalities at $n=1$, where $\nu_i(1)=a^i$.

\begin{theorem}[All rows when the top row is at most one]\label{thm:onepolyrows}
Let $r>1$ and $n\ge1$ with $\beta_r(n)\le1$.  Then
$1\le\nu_1(n)\le\nu_2(n)\le\cdots$, and for every $k\ge1$ and $L=n+k-1$,
\[
  V_r(L,k)=\frac1{\nu_k(n)}=\min_{1\le i\le k}\frac1{\nu_i(n)} .
\]
So the prefix identity~\eqref{eq:prefixid} holds at $(r,n+k-1,k)$
for every $k$, with the minimum at $i=k$.
\end{theorem}

\begin{proof}
By Lemma~\ref{lem:prefixshift} with $m=1$ and $\nu_1(n)=1/\beta_r(n)\ge1$,
$A_{i+1}(p)\ge A_i(p)$ for every $p\in\mathcal Q_n$ and every $i$, and
$\nu_{i+1}(n)\ge\nu_i(n)$.  Given $\varepsilon>0$ choose $p\in\mathcal Q_n$
with $A_k(p)\le\nu_k(n)+\varepsilon$.  Then
$\max_{i\le k}A_i(p)=A_k(p)$, and Theorem~\ref{thm:onepoly} gives
$V_r(L,k)\ge1/(\nu_k(n)+\varepsilon)$.  The upper bound
$V_r(L,k)\le\min_i1/\nu_i(n)\le1/\nu_k(n)$ is \cite[Theorem~3.3]{coefficient-mass-rows}.
\end{proof}

The hypothesis $\beta_r(n)\le1$ holds for $r$ in an interval with left end
$1$, which is $(1,2]$ for $n=1$, shrinks as $n$ grows, and is nonempty for
every $n$ (Lemma~\ref{lem:topnearone}).  Indeed, since
$\nu_1(n+1)\le\nu_1(n)$, we have $\beta_r(n)\ge\beta_r(1)=r-1$.  So the
hypothesis forces $r\le2$, and at $r=2$ it holds only for $n=1$, since
$\beta_2(n)\ge n$ by \cite[Theorem~4.13]{coefficient-mass}.  For each
$q\in\mathcal Q_n$ the sum $\Phi_r(q)$ decreases in $r$, term by term, so
$\nu_1(n)$ does not increase with $r$, and the set
$\{r>1:\beta_r(n)\le1\}$ is an interval with left end $1$.

\begin{lemma}[The top row near the root $1$]\label{lem:topnearone}
Let $n\ge1$, $1<r\le2$ and $C_n=\prod_{i=1}^n(2i+1)$.  Then
$\beta_r(n)\le e^{4n+2}C_n\,(r-1)$.  In particular Theorem~\ref{thm:onepolyrows}
applies whenever $r-1\le e^{-4n-2}/C_n$.
\end{lemma}

\begin{proof}
Let $q\in\mathcal Q_n$ and $m=\lceil1/(r-1)\rceil$, so $1\le(r-1)m\le2$.
For $1\le j\le n$ let $B_j$ be the $m$ integers in $[2jm,2jm+m)$.  Fix
$s_j\in B_j$.  Lagrange interpolation at $s_1,\ldots,s_n$ gives
$1=q(0)=\sum_j\lambda_jq(s_j)$ with
$\lambda_j=\prod_{i\ne j}s_i/(s_i-s_j)$.  Here $s_i<(2i+1)m$ and
$|s_i-s_j|>(2|i-j|-1)m\ge m$ for $i\ne j$, so $|\lambda_j|\le C_n$ and
$1\le C_n\sum_j|q(s_j)|$.  Averaging over the $m^n$ choices of
$(s_1,\ldots,s_n)$ gives $m\le C_n\sum_{s\in B}|q(s)|$, where
$B=B_1\cup\cdots\cup B_n\subseteq[1,(2n+1)m]$.  For $s\le(2n+1)m$,
$r^{-s}\ge e^{-(r-1)(2n+1)m}\ge e^{-4n-2}$, since $\log r\le r-1$.  Hence
$\Phi_r(q)\ge e^{-4n-2}m/C_n\ge e^{-4n-2}/(C_n(r-1))$, and taking the
infimum over $q$ gives $\nu_1(n)\ge e^{-4n-2}/(C_n(r-1))$, which is the
claim since $\beta_r(n)=1/\nu_1(n)$.
\end{proof}

Beyond $\beta_r(n)\le1$, where the argument of Theorem~\ref{thm:onepolyrows}
no longer shows that one polynomial serves all prefixes,
we do not know in general when the prefix identity holds;
Section~\ref{sec:secondrow} reduces the case $k=2$ to finitely many exempted
positions, and Section~\ref{sec:finiterows} every $k$ to finitely many
exempted sets.  Section~\ref{sec:longdiag} shows that a
single polynomial does serve all prefixes once $k$ is large compared with
$n$.

\section{Long diagonals}\label{sec:longdiag}

Theorem~\ref{thm:onepoly} gives the value $1/\nu_k(n)$ as soon
as one polynomial that is nearly optimal for $\nu_k(n)$ has its $k$-th
prefix sum $A_k$ at least as large as the earlier ones.
Theorem~\ref{thm:onepolyrows} obtains this from $\beta_r(n)\le1$.  This
section obtains it for every $1<r<2$, with no condition on $\beta_r(n)$,
once $k$ is large compared with $n$ (Theorem~\ref{thm:longdiag} and
Proposition~\ref{prop:longdiagexplicit}).  So on every diagonal
$L-k=n-1$ the prefix identity can fail in at most finitely many rows
(Corollary~\ref{cor:longdiagfinite}).

The earlier weights $\omega_i$, $i<k$, can exceed $\omega_k$ only on
$[1,2k-3]$, where a suitable polynomial for $\nu_k(n)$ is bounded by $1$;
when this excess is small, the $k$-th prefix sum of that polynomial is the
largest.  By
\cite[Lemma~3.1]{coefficient-mass-rows} the polynomials for
$\nu_k(n)$ may be taken with all zeros at least $k$, and such a polynomial is
bounded by $1$ on $[1,2k-3]$.  There the weights $\omega_i$, $i<k$, exceed
$\omega_k$ by an amount that is exponentially small compared with $\nu_k(n)$
when $r<2$ and $k$ is large compared with $n$.  From $2k-2$ on, $\omega_k$ dominates every $\omega_i$, $i<k$,
with a margin that is a fixed fraction of $\omega_k$ near the mode of
$\omega_k$, where every polynomial in $\mathcal Q_n$ has some mass.

For $k\ge2$ put
\[
  w_k(s)=\omega_k(s)-\omega_{k-1}(s),\qquad
  \tau_k(n)=\inf_{p\in\mathcal Q_n}\sum_{s\ge2k-2}w_k(s)\,|p(s)|,
\]
\[
  N_k(r)=\max_{1\le i<k}\sum_{s=1}^{2k-3}\bigl(\omega_i(s)-\omega_k(s)\bigr)^+,
\]
where $x^+=\max\{x,0\}$ and $\omega_i(s)=0$ for $s<i$.  For $s\ge k$ we have
$\omega_{k-1}(s)/\omega_k(s)=(k-1)/(s-k+1)$, so
\begin{equation}\label{eq:wkratio}
  w_k(s)=\frac{s-2k+2}{s-k+1}\,\omega_k(s)\qquad(s\ge k).
\end{equation}
In particular $w_k\ge0$ on $[2k-2,\infty)$, and $\tau_k(n)\ge0$.

\begin{lemma}[Far zeros keep the last prefix sum largest]\label{lem:farprefix}
Let $r>1$, $k\ge2$ and $n\ge1$.  Let $Y$ be a set of $n-1$ integers, each at
least $k$, and $p=r_Y$.  Then for $1\le i<k$,
\[
  A_k(p)-A_i(p)\ \ge\ \tau_k(n)-N_k(r).
\]
\end{lemma}

\begin{proof}
Let $1\le s\le2k-3$.  Every $y\in Y$ has $2y\ge2k>s$, so $|1-s/y|\le1$ and
$|p(s)|\le1$.  Hence
$(\omega_k(s)-\omega_i(s))|p(s)|\ge-(\omega_i(s)-\omega_k(s))^+$, and the
sum of these terms over $s\le2k-3$ is at least $-N_k(r)$.

Let $s\ge2k-2$.  The coefficients $\binom{s-1}j$ increase in $j$ for
$j\le(s-1)/2$, and $i-1\le k-2\le(s-1)/2$, so $\omega_i(s)\le\omega_{k-1}(s)$
and $(\omega_k(s)-\omega_i(s))|p(s)|\ge w_k(s)|p(s)|$.  Since $p\in\mathcal Q_n$,
the sum of these terms over $s\ge2k-2$ is at least $\tau_k(n)$.
\end{proof}

\begin{theorem}[Long diagonals]\label{thm:longdiag}
Let $r>1$, $n\ge1$, $k\ge2$ and $L=n+k-1$.  If $N_k(r)\le\tau_k(n)$, then
$\nu_i(n)\le\nu_k(n)$ for $1\le i\le k$, and
\[
  V_r(L,k)=\frac1{\nu_k(n)}=\min_{1\le i\le k}\frac1{\nu_i(n)} .
\]
So the prefix identity~\eqref{eq:prefixid} holds at
$(r,L,k)$, with the minimum at $i=k$.
\end{theorem}

\begin{proof}
Let $\varepsilon>0$.  \cite[Lemma~3.1]{coefficient-mass-rows} with $\omega=\omega_k$ gives a
set $Y$ of $n-1$ integers, each at least $k$, such that $p=r_Y$ has
$A_k(p)\le\nu_k(n)+\varepsilon$.  By Lemma~\ref{lem:farprefix} and the
hypothesis, $A_i(p)\le A_k(p)$ for $i<k$.  Hence
$\nu_i(n)\le A_i(p)\le\nu_k(n)+\varepsilon$ for $i\le k$, and
$\max_{i\le k}A_i(p)=A_k(p)$.  Theorem~\ref{thm:onepoly} gives
$V_r(L,k)\ge1/(\nu_k(n)+\varepsilon)$.  Let $\varepsilon\to0$.
\cite[Theorem~3.3]{coefficient-mass-rows} gives
$V_r(L,k)\le\min_i1/\nu_i(n)\le1/\nu_k(n)$.
\end{proof}

The hypothesis compares a quantity that is exponentially small relative to
$a^k$ for $r<2$ with one that is not.  The next proposition makes this explicit.  It uses
the mode of $\omega_k$ and Lagrange interpolation, as in
\cite[Lemma~3.4]{coefficient-mass-rows}.

\begin{proposition}[An explicit form]\label{prop:longdiagexplicit}
Let $1<r<2$, $\theta=4(r-1)/r^2$ and $n\ge1$.  For $k\ge2$ put
$M_k=\lfloor(k-1)r/(r-1)\rfloor+1$ and
\[
  \Xi(k)=\theta^k\,(2r)^{n-1}\binom{M_k+n-1}{n-1}
  \Bigl(\frac{4\sqrt{kr}}{r-1}+1\Bigr).
\]
\begin{enumerate}
\item[(a)] If $\Xi(k)\le3(2-r)^2/r^2$, then $N_k(r)\le\tau_k(n)$, so
Theorem~\ref{thm:longdiag} applies.
\item[(b)] Put $d=\lceil r/(r-1)\rceil$ and
$\gamma_k=\theta\,(1+(n-1)/(M_k+1))^d\sqrt{1+1/k}$.  Then $\Xi(k+1)\le\gamma_k\Xi(k)$,
and $\gamma_k$ does not increase with $k$.  So if $\Xi(k_1)\le3(2-r)^2/r^2$ and
$\gamma_{k_1}\le1$, then (a) applies to every $k\ge k_1$.
\item[(c)] Let $\alpha\ge1$ satisfy
$\rho(\alpha):=\theta\,\bigl(2re(\alpha r/(r-1)+1)\bigr)^{1/\alpha}<1$.  If
$n-1\le k/\alpha$, then
$\Xi(k)\le\rho(\alpha)^k(4\sqrt{kr}/(r-1)+1)$.  Consequently there is $K$ such
that $V_r(n+k-1,k)=1/\nu_k(n)$ for every $n\ge1$ and every
$k\ge\max\{K,\alpha(n-1)\}$.
\end{enumerate}
Since $\bigl(2re(\alpha r/(r-1)+1)\bigr)^{1/\alpha}\to1$ as $\alpha\to\infty$
and $\theta<1$, admissible $\alpha$ exist for every $1<r<2$.
\end{proposition}

\begin{proof}
(a) First, $(\omega_i(s)-\omega_k(s))^+\le\omega_i(s)\le
\binom{s-1}{\lfloor(s-1)/2\rfloor}r^{-s}$, so $N_k(r)\le E_k(r)$, and the
proof of \cite[Corollary~3.5]{coefficient-mass-rows} gives
$E_k(r)\le\frac{r^2}{4(2-r)}(4/r^2)^k=\frac{r^2}{4(2-r)}\,\theta^ka^k$.

Next let $p\in\mathcal Q_n$.  The ratio $\omega_k(s+1)/\omega_k(s)=s/(r(s-k+1))$
is at least $1/r$, and it is at least $1$ exactly when $s\le(k-1)r/(r-1)$.  So
$M_k$ maximizes $\omega_k$: it is a mode of the law $P$ of
\cite[Lemma~3.4]{coefficient-mass-rows}, where $\omega_k=a^kP$.  That lemma gives
$\omega_k(M_k)\ge a^k\cdot3/(4(4\sqrt{kr}/(r-1)+1))$, and
$\omega_k(M_k+l)\ge r^{-l}\omega_k(M_k)$ for $l\ge0$.  Since
$M_k>(k-1)r/(r-1)\ge2k-2$, the nodes $t_l=M_k+l$, $0\le l<n$, lie in
$[2k-2,\infty)$.  By~\eqref{eq:wkratio} the ratio $w_k(s)/\omega_k(s)=1-(k-1)/(s-k+1)$
increases in $s$, and at $s=t_0>(k-1)r/(r-1)$ it exceeds $1-(r-1)=2-r$.  So
$w_k(t_l)\ge(2-r)\,r^{-(n-1)}\omega_k(M_k)$ for every $l$.  Finally
$1=p(0)=\sum_l\ell_lp(t_l)$ with the Lagrange coefficients
$|\ell_l|=\prod_{l'\ne l}t_{l'}/|l-l'|\le\prod_{j=1}^{n-1}(M_k+j)/(l!\,(n-1-l)!)$,
so $\sum_l|\ell_l|\le2^{n-1}\binom{M_k+n-1}{n-1}$ and
$\max_l|p(t_l)|\ge1/(2^{n-1}\binom{M_k+n-1}{n-1})$.  Altogether
\[
  \sum_{s\ge2k-2}w_k(s)|p(s)|\ge\sum_lw_k(t_l)|p(t_l)|
  \ge\frac{3(2-r)\,a^k}{4\,(4\sqrt{kr}/(r-1)+1)\,(2r)^{n-1}\binom{M_k+n-1}{n-1}} .
\]
This lower bound for $\tau_k(n)$ is at least the upper bound for $E_k(r)$
exactly when $\Xi(k)\le3(2-r)^2/r^2$.

(b) We have $M_{k+1}-M_k\le d$, so
$\binom{M_{k+1}+n-1}{n-1}/\binom{M_k+n-1}{n-1}=\prod_{j=M_k+1}^{M_{k+1}}(j+n-1)/j
\le(1+(n-1)/(M_k+1))^d$.  Also
$(4\sqrt{(k+1)r}/(r-1)+1)/(4\sqrt{kr}/(r-1)+1)\le\sqrt{(k+1)/k}$.  This gives
$\Xi(k+1)\le\gamma_k\Xi(k)$.  Since $M_k$ does not decrease, neither does
$1/\gamma_k$, and induction on $k$ gives the rest.

(c) For $n=1$ the claim is $\theta\le\rho(\alpha)$.  Let $m=n-1\ge1$.  Then
$\binom{M_k+m}m\le(e(M_k+m)/m)^m$, and the right side increases with $m$,
since the derivative of $m\log(e(M_k+m)/m)$ is $\log(1+M_k/m)+m/(M_k+m)>0$.
With $m\le k/\alpha$ and $M_k\le(k-1)r/(r-1)+1\le kr/(r-1)$ this gives
$\binom{M_k+m}m\le(e(\alpha r/(r-1)+1))^{k/\alpha}$.  Also
$(2r)^m\le(2r)^{k/\alpha}$, which proves the bound on $\Xi(k)$.  Its right side
tends to $0$ as $k\to\infty$ and decreases once $k\ge1/(2\log(1/\rho(\alpha)))$.
So some $K$ has $\rho(\alpha)^k(4\sqrt{kr}/(r-1)+1)\le3(2-r)^2/r^2$ for all
$k\ge K$, and (a) and Theorem~\ref{thm:longdiag} apply whenever also
$k\ge\alpha(n-1)$.
\end{proof}

\begin{corollary}[Every diagonal is eventually determined]\label{cor:longdiagfinite}
Let $1<r<2$ and $n\ge1$.  Then $V_r(L,L-n+1)=1/\nu_{L-n+1}(n)$ for all but
finitely many $L\ge n$.  Moreover the sequence $\nu_k(n)$ is eventually
nondecreasing in $k$, and $\nu_k(n)=\max_{i\le k}\nu_i(n)$ for all large $k$.
\end{corollary}

\begin{proof}
Apply Proposition~\ref{prop:longdiagexplicit}(c) with $k=L-n+1\ge\max\{K,\alpha(n-1)\}$.
Theorem~\ref{thm:longdiag} also gives $\nu_i(n)\le\nu_k(n)$ for all
$i\le k$ and all such $k$.
\end{proof}

This sharpens \cite[Corollary~3.5]{coefficient-mass-rows}, which gives
$V_r(L,k)=(1+o(1))/\nu_k(n)$ as $k\to\infty$ with $n$ fixed: the factor
$1+o(1)$ is exactly $1$ once $k\ge\max\{K,\alpha(n-1)\}$, where $K$ does not
depend on $n$.
What remains open for $1<r<2$ are the rows with $\beta_r(n)>1$ and $k$
below these thresholds.  The next section treats the first of them,
$k=2$, and Section~\ref{sec:finiterows} reduces each of them to finitely
many exempted sets.

\section{The second row}\label{sec:secondrow}

This section treats the second row, $k=2$, with no condition on
$\beta_r(n)$, through the values $\eta_\sigma(n)=\mu_r(\{\sigma\},n+1)$ of
single exempted positions $\sigma$.  A polynomial vanishing at $\sigma$ is
$(1-s/\sigma)p$ with $p\in\mathcal Q_n$, so
\[
  \eta_\sigma(n)=\mu_r(\{\sigma\},n+1)
  =\inf_{p\in\mathcal Q_n}\sum_{s\ge1}\Bigl|1-\frac s\sigma\Bigr||p(s)|\,r^{-s}
  \qquad(\sigma\in\N_{>0}).
\]
Then $\eta_1(n)=\nu_2(n)$, and \cite[Proposition~2.1(b)]{coefficient-mass-rows} gives
$V_r(n+1,2)=\inf_\sigma1/\eta_\sigma(n)$.  Far positions cost no more than
the limit $\nu_1(n)$, so the second row is the minimum of $1/\nu_1(n)$ and
finitely many position terms (Theorem~\ref{thm:secondrow}).  Each term is
$1/\eta_\sigma(n)=1/\Phi_r(Z)$ for a set $Z\ni\sigma$ of $n$ integers, and
such a $Z$ can be certified exactly
(Lemma~\ref{lem:vertexopt}).  At $r=\frac54$ and $r=\frac43$
this determines rows on both sides of the prefix identity
(Proposition~\ref{prop:secondrowexact}), and at $r=\frac{177}{167}$ it
shows that the identity can fail although $\nu_2(n)>\nu_1(n)$
(Proposition~\ref{prop:secondrowhyp}).  Since $\eta_\sigma(n)$ and
$\nu_i(n)$ decrease with $r$, one certificate covers an interval of roots,
and $49$ of them show that the identity fails in the second row of some
diagonal at every $r\in[\frac{21}{20},\frac{139}{100}]$
(Theorem~\ref{thm:failinterval}).

\begin{lemma}[Dropping a far exempted position]\label{lem:rowmono}
Let $r>1$ and $2\le k\le L$.  Then $V_r(L,k)\le V_r(L-1,k-1)$.  So
$V_r(n+k-1,k)$ does not increase with $k$, for each $n\ge1$.
\end{lemma}

\begin{proof}
Let $S'\subset\N_{>0}$ with $|S'|=k-2$, let $N>\max(S'\cup\{0\})$ and
$S_N=S'\cup\{N\}$.  A real $q$ admissible for $\mu_r(S_N,L)$ factors as
$q=(1-s/N)q_0$ with $q_0$ admissible for $\mu_r(S',L-1)$.  For
$1\le s\le M<N$ we have $1-s/N\ge1-M/N$, so
$\Phi_r(q)\ge(1-M/N)\,\mu_{r,M}(S',L-1)$.  Hence
$\liminf_N\mu_r(S_N,L)\ge\mu_{r,M}(S',L-1)$ for every $M$, and by the
convergence in the proof of \cite[Proposition~2.1(b)]{coefficient-mass-rows} it is at
least $\mu_r(S',L-1)$.  \cite[Proposition~2.1(b)]{coefficient-mass-rows} gives
$V_r(L,k)\le\inf_N1/\mu_r(S_N,L)\le1/\mu_r(S',L-1)$, and the infimum over
$S'$ is $V_r(L-1,k-1)$, again by \cite[Proposition~2.1(b)]{coefficient-mass-rows}.
\end{proof}

With $V_r(n+i-1,i)\le1/\mu_r([1,i-1],n+i-1)=1/\nu_i(n)$ this recovers the
upper bound of \cite[Theorem~3.3]{coefficient-mass-rows}.  It also propagates failures
along a diagonal: if $V_r(n+1,2)<\min_{i\le k}1/\nu_i(n)$, then the prefix
identity fails at $(r,n+k-1,k)$.

For a real polynomial $p$ and real $\sigma>0$ put
\[
  \Psi_r(p)=\sum_{s\ge1}s\,|p(s)|\,r^{-s},\qquad
  R_p(\sigma)=\sum_{s>\sigma}(s-\sigma)\,|p(s)|\,r^{-s}.
\]
Since $|1-s/\sigma|=1-s/\sigma+2(s/\sigma-1)^+$, multiplying by
$|p(s)|r^{-s}$ and summing (all three series converge absolutely) gives
\[
  \Phi_r\bigl((1-s/\sigma)p\bigr)=\Phi_r(p)-\frac{\Psi_r(p)-2R_p(\sigma)}\sigma .
\]
If $p\ne0$, then $\Psi_r(p)\ge\Phi_r(p)>0$ and
$R_p(\sigma)\le\sum_{s>\sigma}s|p(s)|r^{-s}\to0$, so a zero far enough out
lowers $\Phi_r$ strictly.

\begin{lemma}[Certifying a vertex]\label{lem:vertexopt}
Let $r>1$, let $S\subseteq Z$ be finite subsets of $\N_{>0}$ and
$L=|Z|+1$.  For $y\in Z\setminus S$ put $e_y(s)=s\,r_{Z\setminus\{y\}}(s)$
and
\[
  G_y=\sum_{s\ge1,\ s\notin Z}\sgn\bigl(r_Z(s)\bigr)\,e_y(s)\,r^{-s}.
\]
Then $\Phi_r(Z)=\mu_r(S,L)$ if and only if $|G_y|\le|e_y(y)|\,r^{-y}$ for
every $y\in Z\setminus S$.  If $|G_y|<|e_y(y)|\,r^{-y}$ for every
$y\in Z\setminus S$, then $r_Z$ is the unique minimizer: every other real
$q$ with $\deg q\le L-1$, $q(0)=1$ and $q|_S=0$ has
$\Phi_r(q)>\Phi_r(Z)$.  Moreover, for every finite
$S\subset\N_{>0}$ and $L>|S|$ some $Z\supseteq S$ with $|Z|=L-1$ has
$\Phi_r(Z)=\mu_r(S,L)$.
\end{lemma}

\begin{proof}
The polynomials admissible for $\mu_r(S,L)$ form the affine space
$r_Z+E$, where $E$ is the space of real $d$ with $\deg d\le L-1$, $d(0)=0$
and $d|_S=0$, of dimension $L-1-|S|=|Z\setminus S|$.  Each $e_y$ lies in
$E$, and $e_y(y')=0\ne e_y(y)$ for $y\ne y'$ in $Z\setminus S$, so the $e_y$
form a basis and $d=\sum_yc_ye_y$ has $d(y)=c_ye_y(y)$.  The function
$\Phi_r$ is convex and finite on $r_Z+E$, so $r_Z$ is a minimizer if and
only if every one-sided derivative $\Phi_r'(r_Z;d)$ is nonnegative.  For
$s\notin Z$ we have $r_Z(s)\ne0$ and
$(|r_Z(s)+\varepsilon d(s)|-|r_Z(s)|)/\varepsilon\to\sgn(r_Z(s))\,d(s)$ as
$\varepsilon\downarrow0$; for $s\in S$ both terms vanish; for
$s=y\in Z\setminus S$ the quotient is $|d(y)|$.  The quotients are bounded
by $|d(s)|$ and $\sum_s|d(s)|r^{-s}<\infty$, so by dominated convergence
\[
  \Phi_r'(r_Z;d)=\sum_{y\in Z\setminus S}\bigl(c_yG_y+|c_y|\,|e_y(y)|\,r^{-y}\bigr).
\]
This is nonnegative for all $c$ if and only if it is for each $c$
supported on one $y$, with either sign, that is, if and only if
$|G_y|\le|e_y(y)|r^{-y}$ for every $y$.

If every inequality is strict and $d=\sum_yc_ye_y\ne0$, then some
$c_y\ne0$, and each term of $\Phi_r'(r_Z;d)$ is at least
$|c_y|(|e_y(y)|r^{-y}-|G_y|)\ge0$, with strict inequality when $c_y\ne0$;
so $\Phi_r'(r_Z;d)>0$.  By convexity the quotient
$(\Phi_r(r_Z+\varepsilon d)-\Phi_r(r_Z))/\varepsilon$ does not decrease in
$\varepsilon\in(0,1]$, so
$\Phi_r(r_Z+d)-\Phi_r(r_Z)\ge\Phi_r'(r_Z;d)>0$.  Every admissible
$q\ne r_Z$ is $r_Z+d$ with such a $d$.

For the last claim take integers $D_j\to\infty$ with
$D_j\ge L+\max(S\cup\{0\})$ and $W_j$ as in \cite[Lemma~5.1]{coefficient-mass-rows} for
$\mu_{r,D_j}(S,L)$.  Passing to a subsequence, for each $i$ the $i$-th
smallest element of $W_j$ is eventually constant or tends to $\infty$.  A
factor $1-s/w$ with $w\to\infty$ tends to $1$, so $r_{W_j}\to r_Z$
coefficientwise, where $Z\supseteq S$, $|Z|\le L-1$, consists of the
eventually constant elements.  For each $M$,
$\sum_{s\le M}|r_Z(s)|r^{-s}=\lim_j\sum_{s\le M}|r_{W_j}(s)|r^{-s}$, and for
$D_j\ge M$ each partial sum $\sum_{s\le M}|r_{W_j}(s)|r^{-s}$ is at most
$\Phi_{r,D_j}(r_{W_j})=\mu_{r,D_j}(S,L)\le\mu_r(S,L)$.  So
$\Phi_r(Z)\le\mu_r(S,L)$, with equality because $r_Z$ is admissible.  If
$|Z|<L-1$, then $(1-s/\sigma)r_Z$ is admissible for $\mu_r(S,L)$ for every
$\sigma>0$, and for large $\sigma$ the identity before the lemma gives
$\Phi_r((1-s/\sigma)r_Z)<\Phi_r(Z)=\mu_r(S,L)$, a contradiction.  So
$|Z|=L-1$.
\end{proof}

Beyond $\max Z$ the sign of $r_Z$ is $(-1)^{|Z|}$, and
$\sum_{s\ge0}P(s)x^s=\sum_j(\Delta^jP)(0)\,x^j/(1-x)^{j+1}$ for a
polynomial $P$, where $(\Delta P)(s)=P(s+1)-P(s)$ is the forward
difference.  So for rational $r$ both $\Phi_r(Z)$ and the $G_y$ are
rational numbers computable exactly, and for a given $Z$ the test of the
lemma is finite.  By the last claim, every $\mu_r(S,L)$, in particular every
$\eta_\sigma(n)$, is attained at some such $Z$; the lemma does not say how
to find it.


\begin{theorem}[The second row is a finite minimum]\label{thm:secondrow}
Let $r>1$ and $n\ge1$, and put $\eta_\sigma(n)=\mu_r(\{\sigma\},n+1)$ for
$\sigma\in\N_{>0}$.
\begin{enumerate}
\item[(a)] For every $p\in\mathcal Q_n$ and $\sigma\in\N_{>0}$,
\[
  \eta_\sigma(n)\le\Phi_r(p)-\frac{\Psi_r(p)-2R_p(\sigma)}\sigma .
\]
\item[(b)] Some $p_*\in\mathcal Q_n$ has $\Phi_r(p_*)=\nu_1(n)$.  Fix any
such $p_*$, and let $\sigma_*$ be the least $\sigma\in\N_{>0}$ with
$2R_{p_*}(\sigma)\le\Psi_r(p_*)$.  Then $\eta_\sigma(n)\le\nu_1(n)$ for
$\sigma\ge\sigma_*$, with strict
inequality for $\sigma>\sigma_*$, and
\[
  V_r(n+1,2)=\min\Bigl\{\frac1{\nu_1(n)},\
  \min_{1\le\sigma<\sigma_*}\frac1{\eta_\sigma(n)}\Bigr\},
\]
where the inner minimum is omitted if $\sigma_*=1$.
\item[(c)] The prefix identity~\eqref{eq:prefixid} fails at $(r,n+1,2)$
if and only if $\eta_\sigma(n)>\max\{\nu_1(n),\nu_2(n)\}$ for some
$2\le\sigma<\sigma_*$.
\end{enumerate}
Parts (b) and (c) hold for every choice of $p_*$.
\end{theorem}

\begin{proof}
(a) The polynomial $(1-s/\sigma)p$ is admissible for $\eta_\sigma(n)$, and
the identity before Lemma~\ref{lem:vertexopt} gives its value $\Phi_r$.

(b) For $n=1$, $\mathcal Q_1=\{1\}$.  For $n\ge2$, the values
$p(1),\ldots,p(n-1)$ are affine coordinates on $\mathcal Q_n$ and
$\Phi_r(p)\ge r^{-(n-1)}\sum_{s<n}|p(s)|$, so the sublevel sets of $\Phi_r$ are
bounded; $\Phi_r$ is convex and finite on $\mathcal Q_n$, hence continuous,
and a minimizer exists.  Let $p_*$ be any minimizer.  As noted before
Lemma~\ref{lem:vertexopt}, $\Psi_r(p_*)>0$ and $R_{p_*}(\sigma)\to0$, so $\sigma_*$ exists.  Moreover
$R_{p_*}(\sigma)-R_{p_*}(\sigma+1)=\sum_{s>\sigma}|p_*(s)|r^{-s}>0$, since
$p_*\ne0$ has finitely many zeros.  So $\Psi_r(p_*)-2R_{p_*}(\sigma)\ge0$ for
$\sigma\ge\sigma_*$, strictly for $\sigma>\sigma_*$, and (a) with $p=p_*$
gives the first claim.  By \cite[Proposition~2.1(b) and (c)]{coefficient-mass-rows},
$V_r(n+1,2)=\inf_\sigma1/\eta_\sigma(n)$ and
$V_r(n+1,2)\le\beta_r(n)=1/\nu_1(n)$, while the positions
$\sigma\ge\sigma_*$ contribute terms at least $1/\nu_1(n)$.

(c) The prefix identity at $k=2$ reads
$V_r(n+1,2)=\min\{1/\nu_1(n),1/\nu_2(n)\}$.  If $\sigma_*>1$, then
$1/\nu_2(n)=1/\eta_1(n)$ is one of the terms in (b); if $\sigma_*=1$, then
$\nu_2(n)=\eta_1(n)\le\nu_1(n)$ by (b).  So by (b) the identity fails
exactly when some term with $2\le\sigma<\sigma_*$ is smaller than both.
\end{proof}

Part (a) needs no optimality: any $p\in\mathcal Q_n$ and any $\sigma_0$
with $2R_p(\sigma_0)\le\Psi_r(p)$ give $\eta_\sigma(n)\le\Phi_r(p)$ for all
$\sigma\ge\sigma_0$, since $R_p$ decreases.  The next proposition handles
the far positions this way.

\paragraph{Names of zero sets.}
In Propositions~\ref{prop:secondrowexact}, \ref{prop:secondrowhyp}
and~\ref{prop:thirdrowexact} each part fixes a root $r$ and a top-row degree $n$, and a named set of integers
records $n$ and the positions it must contain: $Y_n$ is a set of $n-1$
integers (a candidate for $\nu_1(n)$), $Y'_n$ is a set of $n$ integers
containing $1$, and, for a set $S$ of exempted positions, $Z_S(n)$ is a set
of $n+|S|-1$ integers containing $S$ (a candidate for $\mu_r(S,n+|S|)$),
written $Z_\sigma(n)$ and $Z_{\sigma,\sigma'}(n)$ for $S=\{\sigma\}$ and
$S=\{\sigma,\sigma'\}$.  No name is used at two different roots.  The names
fix only sizes and prescribed elements; every property of a named set is
proved where it is used.

\begin{proposition}[Second rows at the roots $\frac54$ and $\frac43$]\label{prop:secondrowexact}
Let $r=\frac54$ in (a)--(c), and let
\begin{align*}
  Z_2(15)&=\{1,2,5,8,12,17,22,28,35,44,53,65,78,94,115\},\\
  Y_{15}&=\{1,3,6,9,13,18,24,30,38,47,58,71,87,107\},\\
  Y_{13}&=\{1,3,6,10,15,21,27,36,45,57,72,91\}.
\end{align*}
\begin{enumerate}
\item[(a)] $V_r(16,2)=1/\Phi_r(Z_2(15))=1/\nu_3(14)=6.5694\ldots$, while
$\min_{i\le2}1/\nu_i(15)=\beta_r(15)=1/\Phi_r(Y_{15})=7.6492\ldots$.
\item[(b)] $V_r(17,3)\le1/\Phi_r(Z_2(15))<\beta_r(15)=\min_{i\le3}1/\nu_i(15)$.
So the prefix identity fails at $(\frac54,16,2)$ and at $(\frac54,17,3)$.
\item[(c)] $V_r(14,2)=\beta_r(13)=1/\Phi_r(Y_{13})=4.6297\ldots>1$ and
$\nu_2(13)<\nu_1(13)$.  So the prefix identity holds at $(\frac54,14,2)$
with its minimum at $i=1$: the second row is the top row of
$(x-\frac54)^{13}$.
\item[(d)] Let $r=\frac43$,
$Z_3(18)=\{1,2,3,6,8,11,15,18,23,28,34,41,48,57,66,78,91,109\}$ and
$Y_{18}=\{1,2,4,6,9,12,16,20,25,30,36,44,52,61,72,85,103\}$.  Then
$V_r(19,2)=1/\eta_3(18)=1/\Phi_r(Z_3(18))=1/\nu_4(16)=41.38\ldots$, while
$\eta_2(18)<\nu_1(18)$ and
$\min_{i\le2}1/\nu_i(18)=\beta_r(18)=1/\Phi_r(Y_{18})=42.34\ldots$.  So the
prefix identity fails at $(\frac43,19,2)$, and a worst exempted position
is $\sigma=3$, while $\sigma=2$ is not one.
\end{enumerate}
\end{proposition}

\begin{proof}
The proof is computer-assisted: every number below is an exact rational,
computed as described after Lemma~\ref{lem:vertexopt}, and every comparison
is verified by exact rational arithmetic in \texttt{tests/test\_rows.py}
(\texttt{test\_second\_row\_at\_five\_fourths\_fails} for (a) and (b),
\texttt{test\_second\_row\_at\_five\_fourths\_holds} for (c),
\texttt{test\_second\_row\_at\_four\_thirds} for (d)).  No floating
point is used, and the decimals in the statement are truncations of the
exact values, checked by exact comparison with the stated digits.  The
finite computation is this.  For a set $Z$ of integers, $\Phi_r(Z)$ is the
finite sum over $s\le\max Z$ plus the closed form of the series past
$\max Z$; the vertex test of Lemma~\ref{lem:vertexopt} for $Z$ and $S$
evaluates $G_y$, a finite sum plus such a closed form, and
$|e_y(y)|r^{-y}$ for each of the at most $|Z|$ elements $y\in Z\setminus
S$; the tail condition at $\sigma_0$ is the sign of
$2R_p(\sigma_0)-\Psi_r(p)$, evaluated in the same way; and for each
$\sigma<\sigma_0$ the test evaluates $\Phi_r$ on the listed candidate
sets through $\sigma$, at most $n+1$ for each $\sigma$, until one
satisfies the bound.  In (a), (c) and (d) $p=r_{Y_n}$, and $\sigma_0$ is the least
integer $\sigma\ge\max Y_n$ with $2R_p(\sigma)\le\Psi_r(p)$, found by
evaluating this sign at $\sigma=\max Y_n,\max Y_n+1,\ldots$; in all three
parts it is $\max Y_n$ itself.  Any $\sigma_0$ satisfying the tail condition
would do, and none of the proofs uses the least one.  The named sets are the
certificate: with them the statement can be rechecked by any exact
evaluation of $\Phi_r$ and $G_y$.

(a) Lemma~\ref{lem:vertexopt} with $S=\{2\}$ shows
$\eta_2(15)=\Phi_r(Z_2(15))$, with every inequality strict, so
$r_{Z_2(15)}$ is the unique minimizer, and with $S=\varnothing$ that
$\nu_1(15)=\Phi_r(Y_{15})$, which is less than $\Phi_r(Z_2(15))$.  The tail
condition $2R_p(\sigma_0)\le\Psi_r(p)$ holds for $p=r_{Y_{15}}$ at
$\sigma_0=\max Y_{15}=107$, so $\eta_\sigma(15)\le\Phi_r(Y_{15})$ for
$\sigma\ge107$ by the remark after Theorem~\ref{thm:secondrow}.  (Since
$r_{Y_{15}}$ is a minimizer for $\nu_1(15)$, Theorem~\ref{thm:secondrow}(b)
applies with $p_*=r_{Y_{15}}$, and its threshold $\sigma_*$, the least
$\sigma\ge1$ with the tail condition, is $25$, as recorded in
Section~\ref{sec:finiterows} and checked there by exact arithmetic.  The
argument here uses only $\sigma_0=107$, which is larger than necessary,
and not the value of $\sigma_*$.)  For every
$\sigma<107$ other than $2$, one of the following sets $Z\ni\sigma$ of $15$ integers has
$\Phi_r(Z)\le\Phi_r(Z_2(15))$:
$Y'_{15}=\{1,3,5,8,12,17,22,28,36,44,54,65,78,94,115\}$ if
$\sigma\in Y'_{15}$; $\mathrm{ins}(Y_{15},\sigma)$ if $\sigma\notin Y_{15}$;
$Y'_{15}$ with some element replaced by $\sigma$ (all $15$ replacements are
tried) if $\sigma\notin Y'_{15}$.
For each $\sigma$, \texttt{tests/test\_rows.py} generates these candidates
and checks that one of them satisfies the bound.  Hence
$\sup_\sigma\eta_\sigma(15)=\Phi_r(Z_2(15))$, and
\cite[Proposition~2.1(b)]{coefficient-mass-rows} gives $V_r(16,2)=1/\Phi_r(Z_2(15))$.  Since
$[1,2]\subseteq Z_2(15)$,
$\eta_2(15)\le\nu_3(14)=\mu_r([1,2],16)\le\Phi_r(Z_2(15))=\eta_2(15)$.  Finally
$\nu_2(15)\le\Phi_r(Y'_{15})<\nu_1(15)$, as $1\in Y'_{15}$.

(b) Lemma~\ref{lem:rowmono} and (a) give the first inequality.  With (a),
the set $\{1,2,5,8,11,16,21,26,33,41,50,60,71,85,102,123\}\supseteq[1,2]$
gives $\nu_3(15)<\nu_1(15)$, so
$\min_{i\le3}1/\nu_i(15)=1/\nu_1(15)=\beta_r(15)$, which exceeds
$1/\Phi_r(Z_2(15))$ by (a).

(c) Lemma~\ref{lem:vertexopt} gives $\nu_1(13)=\Phi_r(Y_{13})<1$.  The
tail condition holds for $p=r_{Y_{13}}$ at $\sigma_0=\max Y_{13}=91$, and every
$\sigma<91$ lies in a set of $13$ integers with
$\Phi_r\le\Phi_r(Y_{13})$, found among the candidates of (a) built from
$Y_{13}$ and
$Y'_{13}=\{1,3,6,9,14,19,25,33,41,52,64,79,99\}$.  So
$\sup_\sigma\eta_\sigma(13)=\nu_1(13)$, and
\cite[Proposition~2.1(b)]{coefficient-mass-rows} gives $V_r(14,2)=\beta_r(13)$.  Since
$1\in Y'_{13}$, $\nu_2(13)\le\Phi_r(Y'_{13})<\nu_1(13)$.

(d) As in (a): Lemma~\ref{lem:vertexopt} gives $\eta_3(18)=\Phi_r(Z_3(18))$,
with every inequality strict, so $r_{Z_3(18)}$ is the unique minimizer,
and $\nu_1(18)=\Phi_r(Y_{18})<\Phi_r(Z_3(18))$; the tail condition holds for
$p=r_{Y_{18}}$ at $\sigma_0=\max Y_{18}=103$; and every $\sigma<103$ other than $3$
lies in a set of $18$ integers with $\Phi_r\le\Phi_r(Z_3(18))$, found among
the candidates of (a) built from $Y_{18}$ and
$Y'_{18}=\{1,2,4,6,8,11,15,19,23,28,34,41,48,57,66,78,91,109\}$.  Since
$\{1,2\}\subseteq Y'_{18}$, $\eta_2(18)$ and $\nu_2(18)$ are at most
$\Phi_r(Y'_{18})<\nu_1(18)$, and $[1,3]\subseteq Z_3(18)$ gives
$\eta_3(18)=\nu_4(16)=\mu_r([1,3],19)$ as in (a).
\end{proof}

In (a) and (d) the second row equals a prefix value of a longer prefix and
a smaller degree than the prefix identity allows, $1/\nu_3(14)$ and
$1/\nu_4(16)$, although only one position is exempted.  The reason is that
the row is attained at a position $\sigma$ for which the unique minimizer
$r_Z$ of $\eta_\sigma(n)$ has all of $[1,\sigma]$ among its zeros, so that
$\eta_\sigma(n)\le\mu_r([1,\sigma],n+1)\le\Phi_r(Z)=\eta_\sigma(n)$ and the
row is $1/\mu_r([1,\sigma],n+1)=1/\nu_{\sigma+1}(n+1-\sigma)$.  Neither
Theorem~\ref{thm:onepolyrows} nor Theorem~\ref{thm:longdiag} reaches (c),
since both place the minimum at $i=k$.

\paragraph{The hypothesis $\nu_2(n)\ge\nu_1(n)$.}
In Proposition~\ref{prop:secondrowexact}(a) and~(d) the identity fails
with $\nu_2(n)<\nu_1(n)$.  When
$\nu_2(n)>\nu_1(n)$ the position $\sigma=2$ cannot break it:
$\eta_2(n)\le\nu_2(n)$.  Indeed, by Lemma~\ref{lem:vertexopt} some $Y$ with
$|Y|=n-1$ has $\Phi_r(Y)=\nu_1(n)$, and some $Z\ni1$ with $|Z|=n$ has
$\Phi_r(Z)=\nu_2(n)=\mu_r(\{1\},n+1)$.  If $1\in Y$, then $r_Y$ is
admissible for $\mu_r(\{1\},n+1)$ and $\nu_2(n)\le\nu_1(n)$.  So
$Y\subset[2,\infty)$ and $r_Y(2)\ge0$, while
$r_Z(2)=-\prod_{z\in Z\setminus\{1\}}(1-2/z)\le0$, and
\cite[Lemma~4.1(a)]{coefficient-mass-rows} with $T=\varnothing$, $\sigma=2$, $L=n+1$ and
$\varphi=\nu_2(n)$ gives $\eta_2(n)=\mu_r(\{2\},n+1)\le\nu_2(n)$.  The same
argument covers every $\sigma$ at which $r_Y$ and $r_Z$ take opposite
signs, but not every $\sigma$, and the hypothesis $\nu_2(n)>\nu_1(n)$
does not suffice for the identity at $(r,n+1,2)$.

\begin{proposition}[A second-row failure with $\nu_2(n)>\nu_1(n)$]\label{prop:secondrowhyp}
Let $r=\frac{177}{167}$ and
\begin{align*}
  Y_{50}&=\{2,4,7,10,15,20,27,34,42,51,61,71,83,95,109,123,138,155,172,190,\\
  &\qquad210,230,252,275,299,324,350,378,407,437,469,502,537,574,613,654,696,\\
  &\qquad741,789,839,893,950,1011,1076,1148,1226,1314,1416,1543\},\\
  Y'_{50}&=\{1,4,6,10,15,20,26,33,41,50,59,70,81,93,107,121,136,152,169,187,\\
  &\qquad205,225,247,269,292,316,342,369,397,427,458,490,524,560,597,637,678,\\
  &\qquad721,767,815,866,920,978,1039,1105,1177,1256,1345,1448,1576\},\\
  Z_5(50)&=\{2,4,5,10,14,20,26,33,41,50,59,70,81,93,106,121,136,152,168,186,\\
  &\qquad205,225,246,269,292,316,342,369,397,427,458,490,524,560,597,637,678,\\
  &\qquad721,767,815,866,920,978,1039,1105,1177,1256,1345,1448,1576\}.
\end{align*}
Then $\nu_1(50)=\Phi_r(Y_{50})<\nu_2(50)=\Phi_r(Y'_{50})<
\eta_5(50)=\Phi_r(Z_5(50))$, the reciprocals being $4.2262\ldots$,
$4.2221\ldots$ and $4.2112\ldots$.  So
\[
  V_r(51,2)\le\frac1{\eta_5(50)}<\frac1{\nu_2(50)}=\min_{i\le2}\frac1{\nu_i(50)},
\]
and the prefix identity fails at $(\frac{177}{167},51,2)$ although
$\nu_2(50)>\nu_1(50)$.
\end{proposition}

\begin{proof}
The proof is computer-assisted, by exact rational arithmetic as in
Proposition~\ref{prop:secondrowexact}, in
\texttt{test\_second\_row\_hypothesis\_fails} of
\texttt{tests/test\_rows.py}, and again, with the series summed by a
separate method, in \texttt{test\_second\_row\_hypothesis\_independently};
no floating point is used, and the decimals
are truncations of exact values.  Lemma~\ref{lem:vertexopt} with
$S=\varnothing$, $\{1\}$ and $\{5\}$ gives the three equalities, recalling
$\nu_2(50)=\mu_r(\{1\},51)$; for $Z_5(50)$ every inequality of the lemma is
strict, so $r_{Z_5(50)}$ is the unique minimizer for $\eta_5(50)$.  The two
inequalities between the values are exact comparisons.  Finally
\cite[Proposition~2.1(b)]{coefficient-mass-rows} with $S=\{5\}$ gives
$V_r(51,2)\le1/\mu_r(\{5\},51)=1/\eta_5(50)$.
\end{proof}

The unique minimizer for $\eta_5(50)$ vanishes at $2$, $4$ and $5$ but not
at $1$ or $3$.  So the mechanism of Proposition~\ref{prop:secondrowexact}(a)
and~(d), a minimizer for $\eta_\sigma(n)$ vanishing on all of $[1,\sigma]$,
is absent here.  The margins are small: $\nu_2(50)$ exceeds $\nu_1(50)$ by less
than $0.1$ percent, and $\eta_5(50)$ exceeds $\nu_2(50)$ by less than
$0.3$ percent.  We have not determined $V_r(51,2)$.

Floating-point experiments, not used in any proof, suggest that such
failures need $r$ near $1$ and $n$ large.  On the rationals
$r=p/q\in(1,2)$ with $q\le10$ and $n\le14$, and on those in $[\frac54,2)$
with $q\le20$ for every $n$ up to the third consecutive one with
$\nu_2(n)<\nu_1(n)$, every $n$ with $\nu_2(n)\ge\nu_1(n)$ gave
$\max_{\sigma\ge2}\eta_\sigma(n)<0.92\,\nu_2(n)$.  At the root $r_n$ where
$\nu_2(n)=\nu_1(n)$, which decreases from $1.52\ldots$ at $n=4$ to
$1.05\ldots$ at $n=58$, the ratio $\max_{\sigma\ge2}\eta_\sigma(n)/\nu_1(n)$,
computed for $4\le n\le28$ and $n=30,34,\ldots,58$, rises from $0.84$ at
$n=4$ to $0.9998$ at $n=42$, $1.0017$ at $n=46$ and $1.0054$ at $n=58$,
with the maximum at $\sigma=3$ for $n\le12$ and at $\sigma=5$ from $n=13$
on; Proposition~\ref{prop:secondrowhyp} takes $r$ just below $r_{50}$.
We do not know whether a worst exempted position, when one exists, is
always small, or which condition on $r$ and $n$ gives the identity at
$k=2$.

\paragraph{Failures on an interval of roots.}
The failures above are at single rational roots.  Since every $\mu_r(S,L)$
decreases with $r$, a lower bound on $\eta_\sigma(n)$ certified at one root
and upper bounds on $\nu_1(n)$ and $\nu_2(n)$ at a smaller root give a
failure at every root between them (Proposition~\ref{prop:failinterval}), and finitely many such
certificates show that the identity fails somewhere in the second row for
every root in an interval (Theorem~\ref{thm:failinterval}).

\begin{lemma}[Monotone in the root]\label{lem:rootmono}
Let $1<r\le r'$, let $S\subset\N_{>0}$ be finite and $L>|S|$.  Then
$\Phi_{r'}(q)\le\Phi_r(q)$ for every real polynomial $q$, and
$\mu_{r'}(S,L)\le\mu_r(S,L)$.  In particular $\nu_i(n)$ and
$\eta_\sigma(n)$ do not increase with $r$.
\end{lemma}

\begin{proof}
Term by term, $|q(s)|r'^{-s}\le|q(s)|r^{-s}$ for $s\ge1$.  The admissible
set of $\mu_r(S,L)$ does not depend on $r$, so the infima compare in the
same way; $\nu_i(n)=\mu_r([1,i-1],n+i-1)$ and
$\eta_\sigma(n)=\mu_r(\{\sigma\},n+1)$.
\end{proof}

\begin{proposition}[A failure certified on an interval]\label{prop:failinterval}
Let $1<r_0<r_1$, $n\ge2$ and $\sigma\ge2$.  Let $Y$, $Y'$ and $Z$ be sets
of $n-1$, $n$ and $n$ positive integers with $1\in Y'$ and $\sigma\in Z$.
Suppose that $Z$ passes the test of Lemma~\ref{lem:vertexopt} for
$S=\{\sigma\}$ at the root $r_1$, and that
\[
  \Phi_{r_1}(Z)>\max\bigl\{\Phi_{r_0}(Y),\ \Phi_{r_0}(Y')\bigr\}.
\]
Then for every $r\in[r_0,r_1]$, $\eta_\sigma(n)>\max\{\nu_1(n),\nu_2(n)\}$
at the root $r$, so
\[
  V_r(n+1,2)\le\frac1{\eta_\sigma(n)}<\min_{i\le2}\frac1{\nu_i(n)},
\]
and the prefix identity~\eqref{eq:prefixid} fails at $(r,n+1,2)$.
\end{proposition}

\begin{proof}
Write $\eta_\sigma(n;r)$ and $\nu_i(n;r)$ for the values at the root $r$.
Lemma~\ref{lem:vertexopt} with $L=|Z|+1=n+1$ gives
$\eta_\sigma(n;r_1)=\mu_{r_1}(\{\sigma\},n+1)=\Phi_{r_1}(Z)$, and
Lemma~\ref{lem:rootmono} gives $\eta_\sigma(n;r)\ge\Phi_{r_1}(Z)$.  The
polynomial $r_Y$ lies in $\mathcal Q_n$, and $r_{Y'}$ has degree $n$, takes
the value $1$ at $0$ and vanishes at $1$, so it is admissible for
$\nu_2(n;r)=\mu_r(\{1\},n+1)$.  By Lemma~\ref{lem:rootmono},
$\nu_1(n;r)\le\Phi_r(Y)\le\Phi_{r_0}(Y)$ and
$\nu_2(n;r)\le\Phi_r(Y')\le\Phi_{r_0}(Y')$.  So
$\eta_\sigma(n;r)>\max\{\nu_1(n;r),\nu_2(n;r)\}$, and
\cite[Proposition~2.1(b)]{coefficient-mass-rows} with $S=\{\sigma\}$ gives
$V_r(n+1,2)\le1/\eta_\sigma(n;r)$.
\end{proof}

Only $Z$ needs an optimality certificate; $Y$ and $Y'$ only give upper
bounds, and need not be optimal at either root.

\begin{theorem}[Second-row failures on an interval]\label{thm:failinterval}
For every $r\in[\frac{21}{20},\frac{139}{100}]$ there is an $n$ with
$10\le n\le68$ such that $\eta_2(n)>\max\{\nu_1(n),\nu_2(n)\}$ at the root
$r$; so $V_r(n+1,2)\le1/\eta_2(n)<\min_{i\le2}1/\nu_i(n)$, and the prefix
identity~\eqref{eq:prefixid} fails at $(r,n+1,2)$.
\end{theorem}

\begin{proof}
The proof is computer-assisted, by exact rational arithmetic as in
Proposition~\ref{prop:secondrowexact}, in
\texttt{test\_second\_row\_fails\_on\_an\_interval} of
\texttt{tests/test\_rows.py}.  The file \texttt{tests/second\_row\_cover.json}
lists $49$ tuples $(r_0,r_1,n,2,Y,Y',Z)$, with rational endpoints over
denominators dividing $2\cdot10^4$ and $n$ nonincreasing from $68$ to $10$.  The
test checks that each tuple satisfies the hypotheses of
Proposition~\ref{prop:failinterval} with $\sigma=2$: the sizes, $1\in Y'$,
$2\in Z$, the vertex test for $Z$ at $r_1$, and the strict comparison.  It
also checks that the first interval starts at $\frac{21}{20}$, that each
starts at or below the largest right end before it, and that the last one
ends at or above $\frac{139}{100}$.  So the intervals $[r_0,r_1]$ cover
$[\frac{21}{20},\frac{139}{100}]$, and Proposition~\ref{prop:failinterval}
applies at every root there.
\end{proof}

So the failures of Proposition~\ref{prop:secondrowexact} at $\frac54$ and
$\frac43$ are not isolated: the second row misses the prefix identity at
every root of $[\frac{21}{20},\frac{139}{100}]$, each time certified at the
position $2$ and a diagonal $n$ that moves with $r$.  By the argument before
Proposition~\ref{prop:secondrowhyp}, $\eta_2(n)>\nu_2(n)$ forces
$\nu_2(n)\le\nu_1(n)$, so none of these failures has
$\nu_2(n)>\nu_1(n)$.  The certificates were found by
floating-point search, which played no part in the proof.  The same
search, also not used in any proof, finds no second-row failure with
$n\le40$ at the roots $1.400$ and $1.432$, $1.436,\ldots,1.460$,
$1.472,\ldots,1.5$ (step $0.004$), nor with $n\le30$ at $1.6$, $1.7$,
$1.8$ and $1.9$; between these the failures continue in narrow windows
of $r$ for single $n$, such as $n=9$ on $[1.403,1.429]$ and $n=8$ near
$1.466$.  In higher rows the identity fails at every root of
$[\frac{139}{100},\frac{19}{10}]$ (Theorem~\ref{thm:blockinterval} in
Section~\ref{sec:finiterows}).  Whether every $1<r<2$ has a failure in some
row, and whether there is one for every $1<r<\frac{21}{20}$, remain open.
There a failure at $(r,n+k-1,k)$, in any row, needs $\beta_r(n)>1$
(Theorem~\ref{thm:onepolyrows}), so by Lemma~\ref{lem:topnearone} its
diagonal $n$ tends to infinity as $r\to1^+$.

\section{Every row is a finite minimum}\label{sec:finiterows}

Theorem~\ref{thm:secondrow} rests on one fact: a zero placed far enough out
does not increase $\Phi_r$.  For $k\ge3$ several exempted positions can
move out together, and a threshold for one far zero does not serve two.  At
$r=\frac54$ and $p_*=r_{Y_{15}}$ of Proposition~\ref{prop:secondrowexact}
the threshold of Theorem~\ref{thm:secondrow} is $\sigma_*=25$, while the
zeros $25$ and $26$ together raise $\Phi_r(p_*)$ by more than half
(exact rational arithmetic, \texttt{test\_far\_zeros} in
\texttt{tests/test\_rows.py}).  Lemma~\ref{lem:farzeros} gives a threshold
for up to $m$ far zeros at once.  It depends on the exempted positions below
it, and the positions that stay below the thresholds form a finite tree.  So
each row is the minimum of the row before it on its diagonal and finitely
many exempted-set terms (Theorem~\ref{thm:finiterows}).

For integers $m,s\ge1$, real $t>0$ and a real polynomial $q$ put
\[
  \phi_m(s,t)=\begin{cases}-s,&s\le t,\\(s-t)(s/t)^{m-1}-t,&s>t,\end{cases}
  \qquad
  \Delta_{q,m}(t)=\sum_{s\ge1}\phi_m(s,t)\,|q(s)|\,r^{-s}.
\]
Since $\phi_1(s,t)=-s+2(s-t)^+$, we have
$\Delta_{q,1}(\sigma)=2R_q(\sigma)-\Psi_r(q)$.

\begin{lemma}[Several far zeros]\label{lem:farzeros}
Let $r>1$, $m\ge1$, and let $q$ be a nonzero real polynomial.
\begin{enumerate}
\item[(a)] $\Delta_{q,m}(t)$ does not increase in $t>0$,
$\Delta_{q,m}\le\Delta_{q,m+1}$, and $\Delta_{q,m}(t)\to-\Psi_r(q)<0$ as
$t\to\infty$.  So there is a least integer $t_m(q)\ge1$ with
$\Delta_{q,m}(t_m(q))\le0$; moreover $\Delta_{q,m}(t)\le0$ for all real
$t\ge t_m(q)$, and $t_m(q)\le t_{m+1}(q)$.
\item[(b)] If $t>0$, $\Delta_{q,m}(t)\le0$, and $U\subset\N_{>0}$ is finite
with $|U|\le m$ and all elements at least $t$, then
$\Phi_r(q\,r_U)\le\Phi_r(q)$.
\end{enumerate}
For $m=1$, $t_1(p_*)$ is the threshold $\sigma_*$ of
Theorem~\ref{thm:secondrow}.
\end{lemma}

\begin{proof}
(a) Fix $s$.  For $t\ge s$, $\phi_m(s,t)=-s$.  For $0<t<s$, the product
$(s-t)(s/t)^{m-1}$ of a positive decreasing and a positive nonincreasing
function of $t$ decreases, so $\phi_m(s,t)$ decreases in $t$ there, and it
tends to $-s$ as $t\uparrow s$.  So $\phi_m(s,\cdot)$ is nonincreasing.  For
$s>t$, $(s/t)^{m-1}\le(s/t)^m$, so $\phi_m\le\phi_{m+1}$.  For fixed $t$,
$|\phi_m(s,t)|$ is at most a polynomial in $s$, so every series here
converges absolutely, and the first two claims follow term by term.  For
$t\ge1$, $|\phi_m(s,t)|\le s^m+s$, and $\sum_s(s^m+s)|q(s)|r^{-s}<\infty$, so
dominated convergence with $\phi_m(s,t)\to-s$ gives the limit.  Finally
$\Psi_r(q)>0$, since $q$ has finitely many zeros.

(b) Let $U\ne\varnothing$ and $u=\min U\ge t$, so $\Delta_{q,m}(u)\le0$ by
(a).  Every $v\in U$ has $|1-s/v|\le\max\{1,s/v\}\le\max\{1,s/u\}$, so
$|r_U(s)|\le|1-s/u|\max\{1,s/u\}^{m-1}$ for $s\ge1$.  Hence
$u(|r_U(s)|-1)\le\phi_m(s,u)$: for $s\le u$ the left side is at most
$(u-s)-u$, and for $s>u$ at most $(s-u)(s/u)^{m-1}-u$.  Multiplying by
$|q(s)|r^{-s}$ and summing gives
$u\bigl(\Phi_r(q\,r_U)-\Phi_r(q)\bigr)\le\Delta_{q,m}(u)\le0$.  For $m=1$,
$\Delta_{p_*,1}(\sigma)\le0$ is the condition $2R_{p_*}(\sigma)\le\Psi_r(p_*)$
defining $\sigma_*$.
\end{proof}

\begin{theorem}[Every row is a finite minimum]\label{thm:finiterows}
Let $r>1$ and $n,m\ge1$.
\begin{enumerate}
\item[(a)] For each finite $N\subset\N_{>0}$ with $|N|<m$ fix a minimizer
$q_N$ of $\mu_r(N,n+|N|)$ (Lemma~\ref{lem:vertexopt}) and put
$T(N)=t_{m-|N|}(q_N)$.  Let $\mathcal S$ be the set of
$S=\{\sigma_1<\cdots<\sigma_m\}\subset\N_{>0}$ with
$\sigma_{j+1}<T(\{\sigma_1,\ldots,\sigma_j\})$ for $0\le j<m$.  Then
$\mathcal S$ is finite, and
\[
  V_r(n+m,m+1)=\min\Bigl\{V_r(n+m-1,m),\
  \min_{S\in\mathcal S}\frac1{\mu_r(S,n+m)}\Bigr\},
\]
where the inner minimum is omitted if $\mathcal S=\varnothing$.
\item[(b)] Let $B>0$, and let $\mathcal T$ be a set of finite subsets of
$\N_{>0}$ of size less than $m$, with $\varnothing\in\mathcal T$.  To each
$N\in\mathcal T$ attach a real polynomial $q_N$ with $\deg q_N\le n+|N|-1$,
$q_N(0)=1$, $q_N|_N=0$ and $\Phi_r(q_N)\le B$, and an integer $T_N\ge1$ with
$\Delta_{q_N,m-|N|}(T_N)\le0$.  Suppose that $N\cup\{\sigma\}\in\mathcal T$
whenever $N\in\mathcal T$, $|N|\le m-2$ and $\max N<\sigma<T_N$
($\max\varnothing=0$), and that $\mu_r(N\cup\{\sigma\},n+m)\le B$ whenever
$N\in\mathcal T$, $|N|=m-1$ and $\max N<\sigma<T_N$.  Then
$\mu_r(S,n+m)\le B$ for every $S\subset\N_{>0}$ with $|S|=m$, and
$V_r(n+m,m+1)\ge1/B$.
\end{enumerate}
\end{theorem}

\begin{proof}
Both parts rest on one step.  Let $S=\{\sigma_1<\cdots<\sigma_m\}$ and
put $N_j=\{\sigma_1,\ldots,\sigma_j\}$ for $0\le j\le m$.  Fix $j<m$ and
let $q$ be real with $\deg q\le n+j-1$, $q(0)=1$, $q|_{N_j}=0$ and
$\Delta_{q,m-j}(\sigma_{j+1})\le0$.  The set
$U=S\setminus N_j$ has $m-j$ elements, all at least $\sigma_{j+1}$, so
Lemma~\ref{lem:farzeros}(b) gives $\Phi_r(q\,r_U)\le\Phi_r(q)$.  The
polynomial $q\,r_U$ has degree at most $n+m-1$, takes the value $1$ at $0$
and vanishes on $S$, so
\begin{equation}\label{eq:farstep}
  \mu_r(S,n+m)\le\Phi_r(q).
\end{equation}

(a) A set in $\mathcal S$ with given $\sigma_1,\ldots,\sigma_j$ has fewer
than $T(N_j)$ choices of $\sigma_{j+1}$, so $\mathcal S$ is finite by
induction on $j$.  \cite[Proposition~2.1(b)]{coefficient-mass-rows} gives
$V_r(n+m,m+1)\le1/\mu_r(S,n+m)$ for each $S$, and Lemma~\ref{lem:rowmono}
gives $V_r(n+m,m+1)\le V_r(n+m-1,m)$.  Conversely let $|S|=m$ with
$S\notin\mathcal S$.  Then some $j<m$ has $\sigma_{j+1}\ge T(N_j)$, so
$\Delta_{q_{N_j},m-j}(\sigma_{j+1})\le0$ by Lemma~\ref{lem:farzeros}(a),
and~\eqref{eq:farstep} gives $\mu_r(S,n+m)\le\mu_r(N_j,n+j)$.  By
\cite[Proposition~2.1(b)]{coefficient-mass-rows}, $1/\mu_r(N_j,n+j)\ge V_r(n+j,j+1)$ (for
$j=0$ both sides are $\beta_r(n)$), and by Lemma~\ref{lem:rowmono},
$V_r(n+j,j+1)\ge V_r(n+m-1,m)$ since $j\le m-1$.  So every $S$ has
$1/\mu_r(S,n+m)$ at least the right side, and
\cite[Proposition~2.1(b)]{coefficient-mass-rows} gives the reverse inequality.

(b) Let $|S|=m$, and let $j<m$ be least with $\sigma_{j+1}\ge T_{N_j}$, if
any; then $N_0,\ldots,N_j\in\mathcal T$ by induction, since
$N_0=\varnothing\in\mathcal T$, and $N_i\in\mathcal T$,
$\sigma_{i+1}<T_{N_i}$, $i\le m-2$ give $N_{i+1}\in\mathcal T$.  If there
is none, then likewise $N_{m-1}\in\mathcal T$ and $\sigma_m<T_{N_{m-1}}$,
so $\mu_r(S,n+m)\le B$ by the leaf hypothesis.  Otherwise
$\Delta_{q_{N_j},m-j}(\sigma_{j+1})\le\Delta_{q_{N_j},m-j}(T_{N_j})\le0$ by
Lemma~\ref{lem:farzeros}(a), and~\eqref{eq:farstep} gives
$\mu_r(S,n+m)\le\Phi_r(q_{N_j})\le B$.  \cite[Proposition~2.1(b)]{coefficient-mass-rows}
gives $V_r(n+m,m+1)\ge1/B$.
\end{proof}

For $m=1$ the set $\mathcal S$ is $\{\{\sigma\}:\sigma<\sigma_*\}$ and (a) is
Theorem~\ref{thm:secondrow}(b).  Applying (a) to $V_r(n+m-1,m)$, and so on
down to $V_r(n,1)=\beta_r(n)$, writes every row $V_r(n+m,m+1)$ as the
minimum of finitely many values $1/\mu_r(S,n+|S|)$, $|S|\le m$, each equal
to $1/\Phi_r(Z)$ for a set $Z\supseteq S$ of $n+|S|-1$ integers
(Lemma~\ref{lem:vertexopt}).  The sets in $\mathcal S$ depend on the
minimizers $q_N$, and the theorem does not say how to find them; part (b)
needs only upper bounds.  If $V_r(n+m,m+1)<V_r(n+m-1,m)$, the infimum
defining $V_r(n+m,m+1)$ is attained, at some $S\in\mathcal S$.  For rational
$r$, $q=r_Z$ and an integer $t\ge1$, $\Delta_{q,m}(t)$ is a rational number
computable exactly
as $\Phi_r(Z)$ is, since past $\max(Z\cup\{t\})$ its summand is a polynomial
times $r^{-s}$.  So for given zero sets every hypothesis of (b), with
$\mu_r(N\cup\{\sigma\},n+m)\le B$ checked through some $\Phi_r(Z)\le B$,
$Z\supseteq N\cup\{\sigma\}$, $|Z|\le n+m-1$, is a finite exact check.

\begin{proposition}[Third rows at the roots $\frac43$ and $\frac32$]\label{prop:thirdrowexact}
Let $r=\frac43$ in (a) and (b).
\begin{enumerate}
\item[(a)] Let
\begin{align*}
  Y_{12}&=\{1,3,6,9,13,17,23,30,39,50,64\},\\
  Z_2(12)&=\{1,2,5,8,12,16,21,28,35,44,56,71\},\\
  Z_{2,3}(12)&=\{2,3,4,7,10,15,19,25,32,40,49,61,77\}.
\end{align*}
Then $V_r(14,3)=1/\Phi_r(Z_{2,3}(12))=7.7168\ldots$,
$V_r(13,2)=1/\Phi_r(Z_2(12))=1/\nu_3(11)=8.5482\ldots$ and
$\beta_r(12)=1/\Phi_r(Y_{12})=8.8501\ldots$, and
$\min_{i\le k}1/\nu_i(12)=\beta_r(12)$ for $k=2,3,4$.  So
$V_r(14,3)<V_r(13,2)<\beta_r(12)$: the prefix identity fails at
$(\frac43,13,2)$ and at $(\frac43,14,3)$, and the third row lies strictly
below the second on this diagonal.  By Lemma~\ref{lem:rowmono} it also
fails at $(\frac43,15,4)$.  The pair $S=\{2,3\}$ is the only
exempted set with $1/\mu_r(S,14)=V_r(14,3)$.
\item[(b)] Let
\begin{align*}
  Y_{14}&=\{1,3,5,8,11,15,20,26,32,40,50,62,77\},\\
  Z_{2,3}(14)&=\{1,2,3,6,9,13,17,22,27,34,41,50,60,73,89\}.
\end{align*}
Then $V_r(15,2)=\beta_r(14)=1/\Phi_r(Y_{14})=15.0390\ldots=\min_{i\le4}1/\nu_i(14)$
and $V_r(16,3)=1/\Phi_r(Z_{2,3}(14))=1/\nu_4(13)=14.6268\ldots$.  So the
prefix identity holds at $(\frac43,15,2)$, with its minimum at $i=1$, and
fails at $(\frac43,16,3)$ and, by Lemma~\ref{lem:rowmono}, at
$(\frac43,17,4)$.  The pair $S=\{2,3\}$ is the only exempted set
with $1/\mu_r(S,16)=V_r(16,3)$.
\item[(c)] Let $r=\frac32$ and $Y_8=\{1,3,6,9,14,20,28\}$.  Then
$V_r(10,3)=V_r(9,2)=\beta_r(8)=1/\Phi_r(Y_8)=8.0259\ldots=\min_{i\le3}1/\nu_i(8)$,
and $\nu_2(8),\nu_3(8)<\nu_1(8)$.  So the prefix identity holds at
$(\frac32,10,3)$ with its minimum at $i=1$.
\end{enumerate}
\end{proposition}

\begin{proof}
The proof is computer-assisted, by exact rational arithmetic as in
Proposition~\ref{prop:secondrowexact}, in \texttt{tests/test\_rows.py}
(\texttt{test\_third\_row\_below\_second} for (a),
\texttt{test\_third\_row\_fails\_alone} for (b),
\texttt{test\_third\_row\_is\_top\_row} for (c)); no floating point is
used, and the decimals are truncations of exact values.  Each lower bound for a row comes from
Theorem~\ref{thm:finiterows}(b) with $q_\varnothing=r_Y$ for the set $Y$
named in that part, $q_N=r_{K_N}$ for sets $K_N\supseteq N$ of at most
$n+|N|-1$ integers, and $T_N$ the least integer with
$\Delta_{q_N,m-|N|}(T_N)\le0$; a leaf $N\cup\{\sigma\}$ is checked through a
set $Z\supseteq N\cup\{\sigma\}$ of at most $n+m-1$ integers with
$\Phi_r(Z)\le B$.  The test builds the tree from the root
$K_\varnothing=Y$.  For a child $N\cup\{\sigma\}$ of a node $N$ with kept
set $K_N$ it forms these candidates: $K_N$ itself if $\sigma\in K_N$;
otherwise $K_N$ with one of the three elements of $K_N\setminus N$ nearest
to $\sigma$ replaced by $\sigma$, and $K_N\cup\{\sigma\}$; and in either
case the fixed set of $n+|N|$ integers of that part, with each position of
$N\cup\{\sigma\}$ that it misses, taken in increasing order, put in place of
one of the two elements outside $N\cup\{\sigma\}$ nearest to it, in every
combination.  The fixed sets are
\begin{align*}
  \text{(a)}\quad&\{1,3,5,8,12,16,21,28,35,44,56,71\},\
  \{1,3,5,8,11,15,20,26,32,40,50,62,77\},\\
  \text{(b)}\quad&\{1,3,5,7,10,14,18,24,30,37,45,55,67,83\},\
  \{1,3,4,7,10,13,17,22,28,34,42,50,61,73,90\},\\
  \text{(c)}\quad&\{1,3,5,8,12,17,24,33\},\ \{1,3,5,8,11,15,21,28,37\},
\end{align*}
of $n$ and $n+1$ integers, used for children of size one and two.  Among
the candidates with at most $n+|N|$ elements the test keeps one with the
least $\Phi_r$ (on ties, the first formed, and nearest elements are taken
toward the smaller one), unless the child is given a named set.  Each
threshold $T_N$ is found by evaluating the sign of $\Delta_{q_N,m-|N|}$ at
integers (it does not increase, Lemma~\ref{lem:farzeros}(a)), and the test
checks $\Delta_{q_N,m-|N|}(T_N)\le0$ exactly.  For the third rows the trees
have $32$, $35$ and $15$ nodes of size one and $126$, $161$ and $33$ leaves
in (a), (b) and (c); at each node and leaf the test checks $\Phi_r\le B$,
or $\Phi_r<B$ where stated, for the set it keeps.  Theorem~\ref{thm:finiterows}(b)
needs only that each node and each leaf have some set of the allowed size
through it with $\Phi_r\le B$ (or $<B$); which one is kept, and so the
tie-breaking above, changes the thresholds $T_N$ and the tree sizes but not
the validity of the certificate.

(a) Lemma~\ref{lem:vertexopt} with $S=\{2,3\}$, $\{2\}$ and $\varnothing$
gives $\mu_r(\{2,3\},14)=\Phi_r(Z_{2,3}(12))$, $\eta_2(12)=\Phi_r(Z_2(12))$ and
$\nu_1(12)=\Phi_r(Y_{12})$, and $\Phi_r(Y_{12})<\Phi_r(Z_2(12))<\Phi_r(Z_{2,3}(12))$.
For $Z_{2,3}(12)$ every inequality of the lemma is strict, so
$r_{Z_{2,3}(12)}$ is the unique minimizer for $\mu_r(\{2,3\},14)$.
With $n=12$, $m=2$, $B=\Phi_r(Z_{2,3}(12))$, $Y=Y_{12}$ (so
$T_\varnothing=33$) and the pair $\{2,3\}$ given $Z_{2,3}(12)$, the test checks
the hypotheses of Theorem~\ref{thm:finiterows}(b), with $\Phi_r<B$ for every
other set used.  So $\mu_r(S,14)\le B$ for all $|S|=2$, with equality only
at $S=\{2,3\}$: a pair outside the tree has $\mu_r(S,14)\le\Phi_r(q_N)<B$
by~\eqref{eq:farstep}.  \cite[Proposition~2.1(b)]{coefficient-mass-rows} gives
$V_r(14,3)=1/B$.  With $m=1$, $B=\Phi_r(Z_2(12))$ and the position $2$ given
$Z_2(12)$, the same check gives $V_r(13,2)\ge1/\Phi_r(Z_2(12))$, and $S=\{2\}$ gives
equality.  Since $[1,2]\subseteq Z_2(12)$,
$\eta_2(12)\le\nu_3(11)=\mu_r([1,2],13)\le\Phi_r(Z_2(12))=\eta_2(12)$.  Finally
the sets $\{1,3,5,8,12,16,21,28,35,44,56,71\}\ni1$,
$\{1,2,5,7,11,15,20,25,32,40,50,61,77\}\supseteq[1,2]$ and
$\{1,2,3,7,10,14,18,23,29,37,45,55,67,83\}\supseteq[1,3]$ give
$\nu_i(12)<\nu_1(12)$ for $i=2,3,4$, so
$\min_{i\le4}1/\nu_i(12)=1/\nu_1(12)$, and Lemma~\ref{lem:rowmono} gives
$V_r(15,4)\le V_r(14,3)$.

(b) As in (a): Lemma~\ref{lem:vertexopt} gives
$\mu_r(\{2,3\},16)=\Phi_r(Z_{2,3}(14))$, with every inequality strict, so
$r_{Z_{2,3}(14)}$ is the unique minimizer, and $\nu_1(14)=\Phi_r(Y_{14})$, which is
smaller.  Theorem~\ref{thm:finiterows}(b) with $n=14$, $m=2$,
$B=\Phi_r(Z_{2,3}(14))$ and $Y=Y_{14}$ ($T_\varnothing=36$), with $\Phi_r<B$
for every set used except $Z_{2,3}(14)$ at $\{2,3\}$, gives $V_r(16,3)=1/B$,
attained only at $\{2,3\}$.  With $m=1$ and $B=\Phi_r(Y_{14})$ it gives
$V_r(15,2)\ge\beta_r(14)$, and \cite[Proposition~2.1(c)]{coefficient-mass-rows} gives
equality.  Since $[1,3]\subseteq Z_{2,3}(14)$,
$\mu_r(\{2,3\},16)\le\mu_r([1,3],16)=\nu_4(13)\le\Phi_r(Z_{2,3}(14))$, so all
three are equal.  The sets $\{1,3,5,7,10,14,18,24,30,37,45,55,67,83\}\ni1$,
$\{1,2,4,7,10,13,17,22,28,34,42,50,61,73,90\}\supseteq[1,2]$ and
$\{1,2,3,6,9,12,16,21,26,32,39,46,55,66,79,96\}\supseteq[1,3]$ give
$\nu_i(14)<\nu_1(14)$ for $i=2,3,4$, and Lemma~\ref{lem:rowmono} gives
$V_r(17,4)\le V_r(16,3)$.

(c) Lemma~\ref{lem:vertexopt} gives $\nu_1(8)=\Phi_r(Y_8)$.
Theorem~\ref{thm:finiterows}(b) with $n=8$, $m=2$, $B=\Phi_r(Y_8)$, $Y=Y_8$,
and the pair $\{2,3\}$ given $\{2,3,4,7,10,15,20,27,37\}$, gives
$V_r(10,3)\ge\beta_r(8)$.  Lemma~\ref{lem:rowmono} and
\cite[Proposition~2.1(c)]{coefficient-mass-rows} give
$V_r(10,3)\le V_r(9,2)\le\beta_r(8)$.  The sets
$\{1,3,5,8,12,17,24,33\}\ni1$ and $\{1,2,5,7,11,15,21,28,37\}\supseteq[1,2]$
give $\nu_2(8),\nu_3(8)<\nu_1(8)$.
\end{proof}

In (b) the prefix identity fails in the third row although it holds in the
second, so this failure does not come from the second row through
Lemma~\ref{lem:rowmono}.  In (a) the third row is strictly below the
second.  In (b) the zero set $Z_{2,3}(14)$ of the unique minimizer for
$\mu_r(\{2,3\},16)$ contains $[1,3]$, as the zero sets of the unique
minimizers for $\eta_\sigma(n)$ in Proposition~\ref{prop:secondrowexact}(a)
and~(d) contain $[1,\sigma]$; in (a) the zero set $Z_{2,3}(12)$ of the
unique minimizer for $\mu_r(\{2,3\},14)$ omits $1$.  The identities at
$(\frac43,15,2)$ and $(\frac32,10,3)$ are outside
Theorems~\ref{thm:onepolyrows} and~\ref{thm:longdiag}, which give
$\nu_1(n)\le\nu_k(n)$.

\paragraph{Failures in higher rows on an interval of roots.}
Above $r=\frac{139}{100}$ the second-row failures of
Theorem~\ref{thm:failinterval} thin out (see the remark after it), but the
identity keeps failing in higher rows, at the block $S=[2,k]$ of exempted
positions; for $k=2$ this is the position $\sigma=2$ of
Theorem~\ref{thm:failinterval}, and for $k=3$ the pair $\{2,3\}$ of
Proposition~\ref{prop:thirdrowexact}.  The value $\mu_r([2,k],L)$ of the
block is at most the prefix value of a longer prefix and a smaller degree,
and at most $(1+\nu_k(n))/(kr)$ (Lemma~\ref{lem:block}).  As in Theorem~\ref{thm:failinterval}, one
certificate covers an interval of roots (Proposition~\ref{prop:blockinterval}),
and finitely many cover $[\frac{139}{100},\frac{19}{10}]$
(Theorem~\ref{thm:blockinterval}).

\begin{lemma}[The exempted block {$[2,k]$}]\label{lem:block}
Let $r>1$, $n\ge2$, $k\ge2$ and $L=n+k-1$.  Then
\[
  \mu_r([2,k],L)\le\min\Bigl\{\nu_{k+1}(n-1),\ \frac{1+\nu_k(n)}{kr}\Bigr\},
\]
with equality in the first bound if some minimizer for $\mu_r([2,k],L)$
vanishes at $1$.  Consequently, if $\mu_r([2,k],L)>\max_{i\le k}\nu_i(n)$,
then $\nu_{k+1}(n-1)>\max_{i\le k}\nu_i(n)$ and $k<(\beta_r(n)+1)/r$.
\end{lemma}

\begin{proof}
Every polynomial admissible for $\mu_r([1,k],L)=\nu_{k+1}(n-1)$ is
admissible for $\mu_r([2,k],L)$, which gives the first bound; a minimizer
for $\mu_r([2,k],L)$ that vanishes at $1$ is admissible for
$\mu_r([1,k],L)$, which gives the equality.  For the second,
Lemma~\ref{lem:vertexopt} gives a set $Z\supseteq[1,k-1]$ of $n+k-2$
integers with $\Phi_r(Z)=\mu_r([1,k-1],L)=\nu_k(n)$.  Put
$Y=Z\setminus[1,k-1]$ and $p=r_Y$, so $|Y|=n-1$, $Y\subset[k,\infty)$,
$A_k(p)=\Phi_r(Z)=\nu_k(n)$ and $p(-1)=\prod_{y\in Y}(1+1/y)\ge1$.  The
polynomial $q(s)=r_{[2,k]}(s)\,p(s-1)/p(-1)$ has degree $L-1$, $q(0)=1$
and $q|_{[2,k]}=0$.  Here $|r_{[2,k]}(1)|=\prod_{j=2}^k(1-1/j)=1/k$, and
for $s\ge k+1$, $|r_{[2,k]}(s)|=\prod_{j=2}^k(s-j)/j=\binom{s-2}{k-1}/k$,
so that $|r_{[2,k]}(s)|\,r^{-s}=\omega_k(s-1)/(kr)$.  Hence, with $t=s-1$,
\[
  \Phi_r(q)=\frac1{kr\,p(-1)}\Bigl(|p(0)|+\sum_{t\ge k}\omega_k(t)|p(t)|\Bigr)
  =\frac{1+A_k(p)}{kr\,p(-1)}\le\frac{1+\nu_k(n)}{kr}.
\]
Finally let $M=\max_{i\le k}\nu_i(n)$ and $\mu_r([2,k],L)>M$.  The first
bound gives $\nu_{k+1}(n-1)>M$, and the second, with $\nu_k(n)\le M$, gives
$krM<1+M$.  Since $kr>1$ and $M\ge\nu_1(n)=1/\beta_r(n)$, this is
$kr-1<1/M\le\beta_r(n)$.
\end{proof}

So a failure at the block needs a prefix value of a longer prefix and a
smaller degree, $\nu_{k+1}(n-1)$, exceeding every prefix value of the
diagonal, and when a minimizer vanishes at $1$ the block term is that prefix
value, as in Proposition~\ref{prop:secondrowexact}(a) and~(d).  It occurs
only in rows $k<(\beta_r(n)+1)/r$, while
Proposition~\ref{prop:longdiagexplicit}(c) excludes rows
$k\ge\max\{K,\alpha(n-1)\}$.

\begin{proposition}[A failure in row $k$ certified on an interval]\label{prop:blockinterval}
Let $1<r_0<r_1$, $n\ge2$, $k\ge2$ and $a_0=1/(r_0-1)$.  Let
$Z\supseteq[2,k]$ be a set of $n+k-2$ positive integers that passes the
test of Lemma~\ref{lem:vertexopt} for $S=[2,k]$ at the root $r_1$.  Let
$1=j_1<\cdots<j_m\le k$, and for each $l$ let $Y_l\supseteq[1,j_l-1]$ be a
set of $n+j_l-2$ positive integers.  Suppose that for every $i\le k$, with
$l$ the largest index such that $j_l\le i$,
\[
  a_0^{\,i-j_l}\,\Phi_{r_0}(Y_l)<\Phi_{r_1}(Z).
\]
Then for every $r\in[r_0,r_1]$, $\mu_r([2,k],n+k-1)>\max_{i\le k}\nu_i(n)$
at the root $r$, so
\[
  V_r(n+k-1,k)\le\frac1{\mu_r([2,k],n+k-1)}<\min_{i\le k}\frac1{\nu_i(n)},
\]
and the prefix identity~\eqref{eq:prefixid} fails at $(r,n+k-1,k)$.
\end{proposition}

\begin{proof}
Write $\nu_i(n;r)$ for the value at the root $r$, and $a=1/(r-1)\le a_0$.
Lemma~\ref{lem:vertexopt} with $L=|Z|+1=n+k-1$ gives
$\mu_{r_1}([2,k],n+k-1)=\Phi_{r_1}(Z)$, and Lemma~\ref{lem:rootmono} gives
$\mu_r([2,k],n+k-1)\ge\Phi_{r_1}(Z)$.  The polynomial $r_{Y_l}$ has degree
$n+j_l-2$ and vanishes on $[1,j_l-1]$, so it is admissible for
$\nu_{j_l}(n;r)=\mu_r([1,j_l-1],n+j_l-1)$, and Lemma~\ref{lem:rootmono}
gives $\nu_{j_l}(n;r)\le\Phi_{r_0}(Y_l)$.  By~\eqref{eq:prefixup} at the
root $r$, $\nu_i(n;r)\le a^{i-j_l}\nu_{j_l}(n;r)\le
a_0^{\,i-j_l}\Phi_{r_0}(Y_l)$ for $j_l\le i$.  So every $\nu_i(n;r)$,
$i\le k$, is below $\mu_r([2,k],n+k-1)$, and
\cite[Proposition~2.1(b)]{coefficient-mass-rows} with $S=[2,k]$ gives the bound on
$V_r(n+k-1,k)$.
\end{proof}

For $k=2$ and $m=2$ this is Proposition~\ref{prop:failinterval} with
$\sigma=2$.  The anchors $j_l$ save listing a zero set for every prefix
value: a certificate needs one only where the previous anchor, raised by
$a_0$ per step, no longer stays below $\Phi_{r_1}(Z)$.

\begin{theorem}[Failures in higher rows on an interval]\label{thm:blockinterval}
For every $r\in[\frac{139}{100},\frac{19}{10}]$ there are $n$ and $k$ with
$13\le n\le22$ and $9\le k\le384$ such that
$\mu_r([2,k],n+k-1)>\max_{i\le k}\nu_i(n)$ at the root $r$; so
$V_r(n+k-1,k)<\min_{i\le k}1/\nu_i(n)$, and the prefix
identity~\eqref{eq:prefixid} fails at $(r,n+k-1,k)$.  With
Theorem~\ref{thm:failinterval}, for every $r\in[\frac{21}{20},\frac{19}{10}]$
the prefix identity fails at some $(r,L,k)$.
\end{theorem}

\begin{proof}
The proof is computer-assisted, by exact rational arithmetic as in
Proposition~\ref{prop:secondrowexact}, in
\texttt{test\_block\_rows\_fail\_on\_an\_interval} of
\texttt{tests/test\_rows.py}.  The file \texttt{tests/block\_row\_cover.json}
lists $27$ tuples $(r_0,r_1,n,k,Z\setminus[2,k],(j_l,Y_l\setminus[1,j_l-1])_l)$,
with rational endpoints over denominators dividing $2\cdot10^4$.  The test
checks that each tuple satisfies the hypotheses of
Proposition~\ref{prop:blockinterval}: the sizes, the anchors
$1=j_1<\cdots<j_m\le k$, the vertex test for $Z$ at
$r_1$, and the strict comparison for every $i\le k$.  It also checks, as for
Theorem~\ref{thm:failinterval}, that the intervals $[r_0,r_1]$ cover
$[\frac{139}{100},\frac{19}{10}]$, and Proposition~\ref{prop:blockinterval}
applies at every root there.
\end{proof}

In every certificate $1\in Z$, so by Lemma~\ref{lem:block} the value
$\Phi_{r_1}(Z)$ is the prefix value $\nu_{k+1}(n-1)$ at $r_1$: the certified
term comes from a block whose minimizer at $r_1$ vanishes on all of $[1,k]$.  The rows
grow as $r$ approaches $2$, from $k=9$ at $\frac{139}{100}$ to
$k=384$ near $\frac{19}{10}$, while $n$ stays between $13$ and
$22$ and the intervals shrink roughly in proportion to $2-r$.  A
floating-point search, not part of any proof, finds further block
failures at $r=1.9$ for $n=14$ in the rows $210\le k\le254$ and for $n=17$ in
the rows $267\le k\le330$.  We do not know whether the identity fails for
every $r<2$, nor whether the failing rows must tend to infinity as
$r\to2$; at $r=2$ it holds in every row (\cite[Theorem~5.3]{coefficient-mass-rows}).

\paragraph{Data availability.}
No datasets were generated or analyzed in this study.  The exact
computations behind Propositions~\ref{prop:secondrowexact},
\ref{prop:secondrowhyp} and~\ref{prop:thirdrowexact} and
Theorems~\ref{thm:failinterval} and~\ref{thm:blockinterval} are in
\texttt{tests/test\_rows.py}, with the certificates of
Theorems~\ref{thm:failinterval} and~\ref{thm:blockinterval} in
\texttt{tests/second\_row\_cover.json} and
\texttt{tests/block\_row\_cover.json}, of the repository
\href{https://github.com/bangyen/coefficient-mass}{bangyen/coefficient-mass};
they need Python~3.12 or later.

\paragraph{Computer assistance.}
Propositions~\ref{prop:secondrowexact}, \ref{prop:secondrowhyp}
and~\ref{prop:thirdrowexact} and Theorems~\ref{thm:failinterval}
and~\ref{thm:blockinterval} are the only computer-assisted results; every other result is proved by hand.  Their
proofs reduce each claim, through Lemma~\ref{lem:vertexopt},
Theorem~\ref{thm:secondrow}(a), Theorem~\ref{thm:finiterows}(b) and
Propositions~\ref{prop:failinterval} and~\ref{prop:blockinterval}, to
finitely many comparisons of explicit rational numbers: the vertex test of
Lemma~\ref{lem:vertexopt} for each named zero set, the sign of
$\Delta_{q,m}$ at each threshold, the bound $\Phi_r(Z)\le B$ (or $<B$) for
a zero set through each exempted position or set below the thresholds,
the comparisons with the stated decimals, and, for
Theorems~\ref{thm:failinterval} and~\ref{thm:blockinterval}, the
comparisons of Propositions~\ref{prop:failinterval}
and~\ref{prop:blockinterval} for each listed tuple and the order of
their endpoints.  The tests
\texttt{test\_second\_row\_*} (including
\texttt{test\_second\_row\_hypothesis\_fails} and
\texttt{test\_second\_row\_fails\_on\_an\_interval}),
\texttt{test\_third\_row\_*} and
\texttt{test\_block\_rows\_fail\_on\_an\_interval} in
\texttt{tests/test\_rows.py} carry out these comparisons in exact rational
arithmetic (Python's \texttt{fractions} module), evaluating $\Phi_r$, $G_y$
and $\Delta_{q,m}$ through the closed form described after
Lemma~\ref{lem:vertexopt}, with no floating point; these computations are
steps of the proofs.  The zero sets were found by search, and only their
verification is part of the proofs.  The value $\sigma_*=25$ and the
comparison for the zeros $25$ and $26$ quoted in
Section~\ref{sec:finiterows} are computed in the same way
(\texttt{test\_far\_zeros}) and are illustrations, not used in any proof.
The floating-point experiments reported after
Proposition~\ref{prop:secondrowhyp} and after
Theorems~\ref{thm:failinterval} and~\ref{thm:blockinterval} are not part
of any proof or test; they
only guided the search for the zero sets and the intervals.
The remaining tests are checks, not proof steps.
\texttt{test\_prefix\_identity\_fails\_below\_two} builds from
$Z_2(15)$, by complementary slackness, an explicit monic multiple of
$(x-\frac54)^{16}$ of degree $300$ with rational coefficients, verifies the
divisibility by $16$ exact synthetic divisions, and checks that its second
largest nonleading coefficient magnitude is at most $6.5695$, an
independent primal check of the upper bound in
Proposition~\ref{prop:secondrowexact}(a).  Other tests check the closed
form for $\Phi_r$ and the vertex test on known cases.  Each check
is paired with a control that it must reject.  The whole file runs in under a
minute.

\paragraph{Use of AI tools.}
Large language model assistants (Claude, Anthropic) were used in drafting
and revising the text, the proofs and the verification code. The author has checked the mathematics and is responsible
for the content.

\paragraph{Funding and competing interests.}
The author received no specific funding for this work and has no competing
interests to declare.

\begin{thebibliography}{9}

\bibitem{dj}
A.~Dubickas and J.~Jankauskas,
\newblock On the reduced height of a polynomial,
\newblock \emph{Publicationes Mathematicae Debrecen} 71 (2007), 325--348,
\newblock \href{https://doi.org/10.5486/PMD.2007.3728}
{doi:10.5486/PMD.2007.3728}.

\bibitem{coefficient-mass}
B.~Pham,
\newblock Displaced-zero certificates and the coefficient mass of polynomial
multiples,
\newblock preprint, 2026, \url{https://github.com/bangyen/coefficient-mass}.

\bibitem{coefficient-mass-rows}
B.~Pham,
\newblock Every row is a top row: coefficient order statistics of multiples
of a power,
\newblock preprint, 2026, \url{https://github.com/bangyen/coefficient-mass}.

\end{thebibliography}

\end{document}
